\documentclass[12pt,a4paper,oneside]{book}

% ============================================================================
% PACKAGES
% ============================================================================
\usepackage[utf8]{inputenc}
\usepackage[T1]{fontenc}
\usepackage{mathptmx} % Times New Roman font
\usepackage{geometry}
\usepackage{setspace}
\usepackage{graphicx}
\usepackage{listings}
\usepackage{xcolor}
\usepackage{amsmath}
\usepackage{amssymb}
\usepackage{hyperref}
\usepackage{fancyhdr}
\usepackage{titlesec}
\usepackage{tocloft}
\usepackage{tabularx}
\usepackage{booktabs}
\usepackage{multirow}
\usepackage{float}
\usepackage{enumitem}
\usepackage{caption}
\usepackage{subcaption}
\usepackage{pmboxdraw} % Support for Unicode box-drawing characters

% ============================================================================
% PAGE GEOMETRY
% ============================================================================
\geometry{
    a4paper,
    left=1.5in,
    right=1in,
    top=1in,
    bottom=1in
}

% ============================================================================
% LINE SPACING
% ============================================================================
\onehalfspacing

% ============================================================================
% HEADERS AND FOOTERS (IEEE Format)
% ============================================================================
\pagestyle{fancy}
\fancyhf{}
\fancyhead[L]{\leftmark}
\fancyhead[R]{\thepage}
\renewcommand{\headrulewidth}{0.4pt}
\setlength{\headheight}{15pt}

% Remove date from plain page style (affects second page)
\fancypagestyle{plain}{%
    \fancyhf{}%
    \fancyfoot[C]{\thepage}%
    \renewcommand{\headrulewidth}{0pt}%
}

% ============================================================================
% IEEE-STYLE TITLE FORMATTING (Bold)
% ============================================================================
% Part formatting - IEEE style with bold
\titleformat{\part}[display]
{\normalfont\huge\bfseries\centering}
{\partname\ \thepart}
{20pt}
{\Huge\bfseries}

% Chapter formatting - IEEE style with bold
\titleformat{\chapter}[display]
{\normalfont\LARGE\bfseries}
{\chaptertitlename\ \thechapter}
{20pt}
{\Huge\bfseries}

% Section formatting - IEEE style with bold
\titleformat{\section}
{\normalfont\Large\bfseries}
{\thesection}
{1em}
{\Large\bfseries}

% Subsection formatting - IEEE style with bold
\titleformat{\subsection}
{\normalfont\large\bfseries}
{\thesubsection}
{1em}
{\large\bfseries}

% Subsubsection formatting - IEEE style with bold
\titleformat{\subsubsection}
{\normalfont\normalsize\bfseries}
{\thesubsubsection}
{1em}
{\normalsize\bfseries}

% ============================================================================
% CAPTION FORMATTING (IEEE Style - Bold)
% ============================================================================
\captionsetup{
    labelfont=bf,
    textfont=bf,
    justification=centering,
    singlelinecheck=false,
    font=small
}

% ============================================================================
% TABLE OF CONTENTS FORMATTING (Bold)
% ============================================================================
\renewcommand{\cftpartfont}{\bfseries\Large}
\renewcommand{\cftchapfont}{\bfseries}
\renewcommand{\cftsecfont}{\bfseries}

% ============================================================================
% HYPERREF SETUP
% ============================================================================
\hypersetup{
    colorlinks=true,
    linkcolor=blue,
    filecolor=magenta,
    urlcolor=cyan,
    citecolor=green,
    pdftitle={A1M: A 64-bit RISC Architecture with Complete Compiler Toolchain},
    pdfauthor={Shahrear Hossain, Riazul Islam},
}

% ============================================================================
% CODE LISTING SETUP - Professional Black and White Theme (Enhanced/Bold)
% ============================================================================
\definecolor{codebg}{rgb}{0.95,0.95,0.95}        % Slightly darker gray background
\definecolor{codeframe}{rgb}{0.2,0.2,0.2}        % Darker frame for contrast
\definecolor{codekeyword}{rgb}{0.0,0.0,0.0}      % Black keywords (bold)
\definecolor{codecomment}{rgb}{0.25,0.25,0.25}   % Darker gray comments
\definecolor{codestring}{rgb}{0.15,0.15,0.15}    % Darker strings
\definecolor{codenumber}{rgb}{0.3,0.3,0.3}       % Darker line numbers

% Base listing style - Professional black and white with enhanced boldness
\lstdefinestyle{base}{
    backgroundcolor=\color{codebg},
    commentstyle=\color{codecomment}\itshape,
    keywordstyle=\color{codekeyword}\bfseries,
    numberstyle=\tiny\color{codenumber},
    stringstyle=\color{codestring},
    basicstyle=\ttfamily\footnotesize\fontseries{b}\selectfont,
    breakatwhitespace=false,
    breaklines=true,
    captionpos=b,
    keepspaces=true,
    numbers=left,
    numbersep=8pt,
    showspaces=false,
    showstringspaces=false,
    showtabs=false,
    tabsize=4,
    frame=single,
    frameround=tttt,
    rulecolor=\color{codeframe},
    framerule=0.8pt,
    xleftmargin=2em,
    framexleftmargin=1.5em,
    aboveskip=1.5em,
    belowskip=1.5em
}

% Verilog style
\lstdefinestyle{verilog}{
    style=base,
    language=Verilog,
    morekeywords={module,endmodule,input,output,wire,reg,always,posedge,
                  negedge,parameter,localparam,begin,end,if,else,case,
                  endcase,default,assign,initial,for}
}

% Assembly style
\lstdefinestyle{assembly}{
    style=base,
    language=[x86masm]Assembler,
    morekeywords={MOVI,ADD,ADDI,SUB,SUBI,MUL,MULI,DIV,DIVI,MOD,MODI,
                  AND,ANDI,OR,ORI,XOR,XORI,SHL,SHLI,SHR,SHRI,SAR,SARI,
                  LOAD,STORE,LOADB,STOREB,JMP,JZ,JNZ,JGT,JLT,JGE,JLE,
                  CALL,RET,CMP,CMPI,PUSH,POP,SYSCALL,NOP,HALT,MOV},
    comment=[l]{;},
    morecomment=[l]{//}
}

% C style
\lstdefinestyle{c}{
    style=base,
    language=C,
    morekeywords={uint64_t,uint32_t,uint16_t,uint8_t,int64_t,int32_t,
                  int16_t,int8_t,bool,size_t,NULL}
}

% Algo style
\lstdefinestyle{algo}{
    style=base,
    morekeywords={fun,return,if,else,while,for,int,bool,char,long,
                  true,false,print,print_int,print_char,read_int,read_char},
    comment=[l]{//},
    morecomment=[s]{/*}{*/},
    morestring=[b]",
    morestring=[b]'
}

% Bash/Shell style
\lstdefinestyle{bash}{
    style=base,
    language=bash,
    morekeywords={echo,cd,ls,make,gcc,./,sudo,install}
}

% ============================================================================
% CHAPTER AND SECTION FORMATTING
% ============================================================================
\titleformat{\chapter}[display]
{\normalfont\huge\bfseries}{\chaptertitlename\ \thechapter}{20pt}{\Huge}
\titlespacing*{\chapter}{0pt}{0pt}{40pt}

\titleformat{\section}
{\normalfont\Large\bfseries}{\thesection}{1em}{}
\titlespacing*{\section}{0pt}{3.5ex plus 1ex minus .2ex}{2.3ex plus .2ex}

\titleformat{\subsection}
{\normalfont\large\bfseries}{\thesubsection}{1em}{}
\titlespacing*{\subsection}{0pt}{3.25ex plus 1ex minus .2ex}{1.5ex plus .2ex}

% ============================================================================
% CUSTOM COMMANDS
% ============================================================================
\newcommand{\code}[1]{\texttt{#1}}
\newcommand{\keyword}[1]{\textbf{\code{#1}}}
\newcommand{\register}[1]{\texttt{R#1}}
\newcommand{\instruction}[1]{\texttt{#1}}
\newcommand{\hex}[1]{\texttt{0x#1}}

% Professional box for important notes
\newcommand{\notebox}[2]{
\begin{center}
\fbox{\parbox{0.9\textwidth}{
\textbf{#1:} #2
}}
\end{center}
}

% ============================================================================
% BEGIN DOCUMENT
% ============================================================================
\begin{document}

% ============================================================================
% TITLE PAGE
% ============================================================================
\begin{titlepage}
    \centering
    \vspace*{1cm}
    
    {\LARGE\bfseries International Islamic University Chittagong\par}
    \vspace{0.5cm}
    {\large Department of Electrical and Electronics Engineering\par}
    \vspace{2cm}
    
    {\Huge\bfseries A1M: A 64-bit RISC Architecture\\[0.5cm]
    with Complete Compiler Toolchain\par}
    \vspace{1cm}
    
    {\Large\itshape A Comprehensive Study of Computer Architecture,\\
    Assembly Language, and Compiler Design\par}
    \vspace{2cm}
    
    {\Large\bfseries Submitted By:\par}
    \vspace{0.5cm}
    {\large Shahrear Hossain Shawon\par}
    \vspace{2cm}
    
    {\Large\bfseries Supervised By:\par}
    \vspace{0.5cm}
    {\large Lecturer Riazul Islam\par}
    {\large Department of Electrical and Electronics Engineering\par}
    {\large International Islamic University Chittagong\par}
    \vspace{2cm}
    
    \vfill
    
    {\small This work is submitted in partial fulfillment of the requirements\\
    for research in Computer Architecture and Compiler Design\par}
\end{titlepage}

% ============================================================================
% DECLARATION
% ============================================================================
\chapter*{Declaration}
\addcontentsline{toc}{chapter}{Declaration}

I hereby declare that the work presented in this thesis titled \textbf{``A1M: A 64-bit RISC Architecture with Complete Compiler Toolchain''} is the result of my own research and has not been submitted elsewhere for any degree or diploma.

All sources of information have been duly acknowledged.

\vspace{2cm}

\noindent
\begin{tabular}{@{}ll}
\textbf{Shahrear Hossain Shawon} & Date: \rule{3cm}{0.4pt} \\[2cm]
\textbf{Supervisor:} & \\
\textbf{Lecturer Riazul Islam} & Date: \rule{3cm}{0.4pt} \\
Department of Electrical and Electronics Engineering & \\
International Islamic University Chittagong & \\
\end{tabular}

% ============================================================================
% ABSTRACT
% ============================================================================
\chapter*{Abstract}
\addcontentsline{toc}{chapter}{Abstract}

This thesis presents A1M, a comprehensive 64-bit RISC architecture with a complete compiler toolchain, encompassing the entire computer science stack from hardware design in Verilog HDL to the Algo high-level programming language. The A1M architecture features a 32-register design, 5-stage pipeline, and over 60 instructions supporting arithmetic, logical, memory, and control flow operations, with a complete software toolchain including a two-pass assembler, cycle-accurate emulator, and full compiler with lexer, parser, semantic analyzer, and code generator, all integrated with VS Code for syntax highlighting and debugging support. This work serves as both a functional computer system and an educational resource demonstrating fundamental concepts in computer architecture, assembly language programming, and compiler construction, with the entire system being open-source and designed for educational purposes.

\textbf{Keywords:} RISC Architecture, Compiler Design, HDL, Assembly Language, Computer Architecture, Educational Computing

\vspace{1cm}

\noindent
\textbf{Supervisor:} Lecturer Riazul Islam\\
Department of Electrical and Electronics Engineering\\
International Islamic University Chittagong

% ============================================================================
% ACKNOWLEDGMENTS
% ============================================================================
\chapter*{Acknowledgments}
\addcontentsline{toc}{chapter}{Acknowledgments}

I would like to express my sincere gratitude to \textbf{Lecturer Riazul Islam}, my supervisor, for his invaluable guidance, continuous support, and expert advice throughout the development of this work; to the Department of Electrical and Electronics Engineering at International Islamic University Chittagong for providing the research environment and resources necessary for this project; to the faculty members for their valuable feedback and encouragement; to the open-source community for providing excellent tools and documentation, particularly the developers of GCC, LLVM, and various HDL synthesis tools; and to my family and friends for their unwavering support throughout this journey.

Finally, I dedicate this work to all students and researchers interested in understanding the fundamental principles of computer systems, from transistors to compilers.

\vspace{1cm}

\noindent
Shahrear Hossain Shawon\\
Department of Electrical and Electronics Engineering\\
International Islamic University Chittagong

% ============================================================================
% TABLE OF CONTENTS
% ============================================================================
\setcounter{tocdepth}{3}  % Show subsections in TOC
\setcounter{secnumdepth}{3}  % Number subsections
\tableofcontents
\newpage

\listoffigures
\newpage

\listoftables
\newpage

\lstlistoflistings
\newpage

% ============================================================================
% PART I: HARDWARE ARCHITECTURE
% ============================================================================
\part{Hardware Architecture}

% ============================================================================
% CHAPTER 1: INTRODUCTION
% ============================================================================
\chapter{Introduction to A1M Architecture}

\section{Overview}

The A1M architecture represents a complete educational computer system designed to demonstrate fundamental principles of computer architecture, assembly language programming, and compiler design. This 64-bit RISC architecture provides a clean, understandable foundation for learning how computers work at every level of abstraction.

\subsection{Design Philosophy}

The A1M architecture adheres to RISC (Reduced Instruction Set Computer) principles:

\begin{itemize}
    \item \textbf{Simplicity:} Fixed-length 32-bit instructions
    \item \textbf{Regularity:} Consistent instruction formats
    \item \textbf{Orthogonality:} Independent instruction components
    \item \textbf{Performance:} Pipeline-friendly design
    \item \textbf{Extensibility:} Room for future enhancements
\end{itemize}

\subsection{Key Features}

\begin{itemize}
    \item 64-bit data path and registers
    \item 32 general-purpose registers (R0-R31)
    \item 32-bit fixed-length instruction encoding
    \item Harvard architecture (separate instruction and data memory)
    \item 5-stage pipeline (IF, ID, EX, MEM, WB)
    \item Over 60 instructions covering all essential operations
    \item Support for 64-bit addressing (16 exabytes theoretical)
\end{itemize}

\subsection{RISC vs. CISC}

Table \ref{tab:risc-cisc} compares RISC and CISC architectures:

\begin{table}[H]
\centering
\caption{RISC vs. CISC Comparison}
\label{tab:risc-cisc}
\begin{tabularx}{\textwidth}{|X|X|}
\hline
\textbf{RISC (A1M)} & \textbf{CISC (x86)} \\
\hline
Fixed-length instructions (32-bit) & Variable-length instructions (1-15 bytes) \\
\hline
Simple addressing modes & Complex addressing modes \\
\hline
Load/Store architecture & Register-memory operations \\
\hline
Large register file (32 registers) & Small register file (8-16 registers) \\
\hline
Hardwired control & Microcode control \\
\hline
Pipeline-friendly & Pipeline stalls common \\
\hline
\end{tabularx}
\end{table}

\section{System Architecture}

Figure \ref{fig:system-overview} shows the high-level system architecture:

\begin{figure}[H]
\centering
\begin{verbatim}
┌─────────────────────────────────────────────────────────┐
│                    A1M CPU System                       │
├─────────────────────────────────────────────────────────┤
│                                                         │
│  ┌───────────┐    ┌──────────┐    ┌──────────────┐      │
│  │           │    │          │    │              │      │
│  │ Registers │<---│   ALU    │<---│ Control Unit │      │
│  │  (R0-R31) │    │          │    │              │      │
│  │           │    └──────────┘    └──────────────┘      │
│  └───────────┘                                          │
│       ^                                                 │
│       │                                                 │
│       v                                                 │
│  ┌─────────────────────────────────────────────────┐    │
│  │              Memory Subsystem                   │    │
│  │  ┌─────────────────┐  ┌────────────────────┐    │    │
│  │  │ Instruction Mem │  │    Data Memory     │    │    │
│  │  │   (I-Cache)     │  │    (D-Cache)       │    │    │
│  │  └─────────────────┘  └────────────────────┘    │    │
│  └─────────────────────────────────────────────────┘    │
│                                                         │
└─────────────────────────────────────────────────────────┘
\end{verbatim}
\caption{A1M System Overview}
\label{fig:system-overview}
\end{figure}

\section{Thesis Organization}

This thesis is organized into eight parts:

\begin{description}
    \item[Part I:] Hardware Architecture --- Registers, ISA, ALU, Control Unit, Memory, Pipeline
    \item[Part II:] Assembly Language --- Assembly fundamentals, data movement, arithmetic, control flow, system calls, assembler implementation
    \item[Part III:] Emulator --- Emulator architecture and system call implementation
    \item[Part IV:] Algo Language --- High-level language introduction, expressions, control structures, functions, arrays
    \item[Part V:] Compiler Implementation --- Lexical analysis, parsing, semantic analysis, code generation
    \item[Part VI:] Compilation Process --- Build system, optimization techniques
    \item[Part VII:] Practical Examples --- Complete programs, debugging, performance analysis
    \item[Part VIII:] Testing and Verification --- Test suite, future enhancements
\end{description}

% ============================================================================
% CHAPTER 2: REGISTER ARCHITECTURE
% ============================================================================
\chapter{Register Architecture}

\section{Register File Organization}

The A1M architecture implements a 32-register file, providing ample storage for temporary values and reducing memory traffic. Each register is 64 bits wide, supporting full 64-bit integer operations.

\subsection{Register Naming Convention}

Table \ref{tab:registers} shows the complete register set:

\begin{table}[H]
\centering
\caption{A1M Register Set}
\label{tab:registers}
\begin{tabularx}{\textwidth}{|l|l|X|}
\hline
\textbf{Register} & \textbf{ABI Name} & \textbf{Purpose} \\
\hline
R0 & zero & Hardwired to 0 \\
R1 & ret/arg1 & Return value / First syscall argument \\
R2 & arg2 & Second syscall argument \\
R3 & sp & Stack pointer \\
R4-R11 & a0-a7 & Function arguments \\
R12-R30 & t0-t18 & Temporary registers \\
R31 & ra & Return address \\
\hline
\end{tabularx}
\end{table}

\subsection{Special Registers}

\begin{itemize}
    \item \textbf{R0 (zero):} Always contains the value 0. Writes to R0 are ignored. This simplifies many operations and provides a source of constant zero.
    \item \textbf{R3 (sp):} Stack pointer, points to the top of the stack. Grows downward (toward lower addresses).
    \item \textbf{R31 (ra):} Return address, automatically set by the CALL instruction. Contains the address to return to after a function completes.
\end{itemize}

\section{Verilog Implementation}

The register file is implemented in Verilog with dual read ports and one write port:

\begin{lstlisting}[style=verilog, caption={Register File Implementation}, label={lst:regfile}]
module register_file (
    input wire clk,
    input wire reset,
    
    // Read ports
    input wire [4:0] read_addr1,
    input wire [4:0] read_addr2,
    output reg [63:0] read_data1,
    output reg [63:0] read_data2,
    
    // Write port
    input wire write_enable,
    input wire [4:0] write_addr,
    input wire [63:0] write_data
);

// Register array
reg [63:0] registers [0:31];

integer i;

// Initialize registers
initial begin
    for (i = 0; i < 32; i = i + 1) begin
        registers[i] = 64'h0;
    end
end

// Asynchronous read (combinational)
always @(*) begin
    if (read_addr1 == 5'h0)
        read_data1 = 64'h0;  // R0 always zero
    else
        read_data1 = registers[read_addr1];
    
    if (read_addr2 == 5'h0)
        read_data2 = 64'h0;  // R0 always zero
    else
        read_data2 = registers[read_addr2];
end

// Synchronous write (on rising edge)
always @(posedge clk) begin
    if (reset) begin
        for (i = 0; i < 32; i = i + 1) begin
            registers[i] <= 64'h0;
        end
    end else if (write_enable && write_addr != 5'h0) begin
        // Don't write to R0
        registers[write_addr] <= write_data;
    end
end

endmodule
\end{lstlisting}

\section{Register Usage Convention}

\subsection{Caller-Saved vs. Callee-Saved}

\begin{itemize}
    \item \textbf{Caller-saved:} R1-R2, R4-R11, R12-R30, R31 --- The calling function must save these if their values are needed after a function call
    \item \textbf{Callee-saved:} R3 --- Called functions must preserve these registers
\end{itemize}

\subsection{Function Call Example}

\begin{lstlisting}[style=assembly, caption={Register Usage in Function Calls}]
; Caller
caller:
    PUSH R31            ; Save return address
    PUSH R4             ; Save argument registers if needed
    
    MOVI R4, 10         ; First argument
    MOVI R5, 20         ; Second argument
    CALL add_function   ; Call function
    
    MOV R10, R1         ; Save return value
    POP R4              ; Restore saved registers
    POP R31
    RET

; Callee
add_function:
    PUSH R31            ; Save return address
    ADD R1, R4, R5      ; Compute result (R1 = R4 + R5)
    POP R31             ; Restore return address
    RET
\end{lstlisting}

% Continue with Chapter 3...

% ============================================================================
% CHAPTER 3: INSTRUCTION SET ARCHITECTURE
% ============================================================================
\chapter{Instruction Set Architecture}

\section{Instruction Formats}

The A1M ISA uses 32-bit fixed-length instructions in four formats:

\subsection{R-Type (Register)}

Used for register-to-register operations:

\begin{verbatim}
 31    24 23  19 18  14 13   9 8       0
┌────────┬──────┬──────┬──────┬─────────┐
│ Opcode │  RD  │ RS1  │ RS2  │ Unused  │
└────────┴──────┴──────┴──────┴─────────┘
  8 bits  5 bits 5 bits 5 bits  9 bits
\end{verbatim}

\begin{itemize}
    \item \textbf{Opcode:} 8-bit operation code
    \item \textbf{RD:} 5-bit destination register
    \item \textbf{RS1:} 5-bit source register 1
    \item \textbf{RS2:} 5-bit source register 2
\end{itemize}

\subsection{I-Type (Immediate)}

Used for operations with immediate values:

\begin{verbatim}
 31    24 23  19 18  14 13           0
┌────────┬──────┬──────┬──────────────┐
│ Opcode │  RD  │ RS1  │  Immediate   │
└────────┴──────┴──────┴──────────────┘
  8 bits  5 bits 5 bits   14 bits
\end{verbatim}

\begin{itemize}
    \item \textbf{Immediate:} 14-bit signed immediate value (sign-extended to 64 bits)
    \item Range: $-8192$ to $+8191$
\end{itemize}

\subsection{J-Type (Jump)}

Used for jump and branch instructions:

\begin{verbatim}
 31    24 23                      0
┌────────┬─────────────────────────┐
│ Opcode │    Jump Address         │
└────────┴─────────────────────────┘
  8 bits        24 bits
\end{verbatim}

\begin{itemize}
    \item \textbf{Jump Address:} 24-bit target address
    \item Address range: $0$ to $16,777,215$ (16 MB)
\end{itemize}

\subsection{M-Type (Memory)}

Used for memory operations (overlaps with I-Type):

\begin{verbatim}
 31    24 23  19 18  14 13           0
┌────────┬──────┬──────┬──────────────┐
│ Opcode │  RD  │ Base │    Offset    │
└────────┴──────┴──────┴──────────────┘
  8 bits  5 bits 5 bits   14 bits
\end{verbatim}

\section{Opcode Assignments}

Table \ref{tab:opcodes} shows the complete opcode assignment:

\begin{table}[H]
\centering
\caption{A1M Opcode Assignment}
\label{tab:opcodes}
\small
\begin{tabularx}{\textwidth}{|l|l|l|X|}
\hline
\textbf{Code} & \textbf{Mnemonic} & \textbf{Format} & \textbf{Operation} \\
\hline
0x00 & NOP & - & No operation \\
0x01 & HALT & - & Halt execution \\
\hline
0x02 & ADD & R & RD = RS1 + RS2 \\
0x03 & ADDI & I & RD = RS1 + IMM \\
0x04 & SUB & R & RD = RS1 - RS2 \\
0x05 & SUBI & I & RD = RS1 - IMM \\
0x06 & MUL & R & RD = RS1 $\times$ RS2 \\
0x07 & MULI & I & RD = RS1 $\times$ IMM \\
0x08 & DIV & R & RD = RS1 / RS2 \\
0x09 & DIVI & I & RD = RS1 / IMM \\
0x0A & MOD & R & RD = RS1 \% RS2 \\
0x0B & MODI & I & RD = RS1 \% IMM \\
\hline
0x0C & AND & R & RD = RS1 \& RS2 \\
0x0D & ANDI & I & RD = RS1 \& IMM \\
0x0E & OR & R & RD = RS1 | RS2 \\
0x0F & ORI & I & RD = RS1 | IMM \\
0x10 & XOR & R & RD = RS1 $\oplus$ RS2 \\
0x11 & XORI & I & RD = RS1 $\oplus$ IMM \\
0x12 & SHL & R & RD = RS1 $<<$ RS2 \\
0x13 & SHLI & I & RD = RS1 $<<$ IMM \\
0x14 & SHR & R & RD = RS1 $>>$ RS2 (logical) \\
0x15 & SHRI & I & RD = RS1 $>>$ IMM (logical) \\
0x16 & SAR & R & RD = RS1 $>>$ RS2 (arithmetic) \\
0x17 & SARI & I & RD = RS1 $>>$ IMM (arithmetic) \\
\hline
0x20 & LOAD & M & RD = Memory[RS1 + Offset] \\
0x21 & STORE & M & Memory[RS1 + Offset] = RD \\
0x22 & LOADB & M & RD = (byte)Memory[RS1 + Offset] \\
0x23 & STOREB & M & Memory[RS1 + Offset] = (byte)RD \\
\hline
0x30 & JMP & J & PC = Address \\
0x31 & JZ & J & if (ZF) PC = Address \\
0x32 & JNZ & J & if (!ZF) PC = Address \\
0x33 & JGT & J & if (!ZF \&\& !NF) PC = Address \\
0x34 & JLT & J & if (NF) PC = Address \\
0x39 & CALL & J & R31 = PC + 4; PC = Address \\
0x3A & RET & - & PC = R31 \\
\hline
0x40 & CMP & R & Set flags: RS1 - RS2 \\
0x41 & CMPI & I & Set flags: RS1 - IMM \\
0x42 & PUSH & R & SP -= 8; Memory[SP] = RS \\
0x43 & POP & R & RD = Memory[SP]; SP += 8 \\
\hline
0x50 & SYSCALL & I & System call (IMM = number) \\
\hline
\end{tabularx}
\end{table}

\section{Instruction Encoding Examples}

\subsection{Example 1: ADD R1, R2, R3}

\begin{verbatim}
Instruction: ADD R1, R2, R3  (R1 = R2 + R3)
Format: R-Type

Binary encoding:
 Opcode    RD     RS1    RS2   Unused
 00000010  00001  00010  00011 000000000
 
Hexadecimal: 0x02 11 23 00 = 0x02112300
\end{verbatim}

\subsection{Example 2: ADDI R4, R5, 100}

\begin{verbatim}
Instruction: ADDI R4, R5, 100  (R4 = R5 + 100)
Format: I-Type

Binary encoding:
 Opcode    RD     RS1    Immediate
 00000011  00100  00101  00000001100100
 
Hexadecimal: 0x03 14 28 64 = 0x03142864
\end{verbatim}

% ============================================================================
% CHAPTER 4: ARITHMETIC LOGIC UNIT (ALU)
% ============================================================================
\chapter{Arithmetic Logic Unit (ALU)}

\section{ALU Overview}

The Arithmetic Logic Unit (ALU) is the computational heart of the A1M processor. It performs all arithmetic and logical operations on 64-bit operands.

\subsection{Supported Operations}

The A1M ALU supports the following operations:

\begin{itemize}
    \item \textbf{Arithmetic:} ADD, SUB, MUL, DIV, MOD
    \item \textbf{Logical:} AND, OR, XOR, NOT
    \item \textbf{Shift:} SHL (left), SHR (logical right), SAR (arithmetic right)
    \item \textbf{Comparison:} CMP (sets flags)
\end{itemize}

\subsection{ALU Control Signals}

Table \ref{tab:alu-ops} shows the ALU operation encoding:

\begin{table}[H]
\centering
\caption{ALU Operation Codes}
\label{tab:alu-ops}
\begin{tabularx}{\textwidth}{|l|l|X|}
\hline
\textbf{Code} & \textbf{Operation} & \textbf{Description} \\
\hline
4'b0000 & ADD & Addition: result = A + B \\
4'b0001 & SUB & Subtraction: result = A - B \\
4'b0010 & MUL & Multiplication: result = A $\times$ B \\
4'b0011 & DIV & Division: result = A / B \\
4'b0100 & MOD & Modulo: result = A \% B \\
4'b0101 & AND & Bitwise AND: result = A \& B \\
4'b0110 & OR & Bitwise OR: result = A | B \\
4'b0111 & XOR & Bitwise XOR: result = A $\oplus$ B \\
4'b1000 & SHL & Shift left: result = A $<<$ B \\
4'b1001 & SHR & Shift right (logical): result = A $>>$ B \\
4'b1010 & SAR & Shift right (arithmetic) \\
4'b1011 & CMP & Compare: set flags based on A - B \\
\hline
\end{tabularx}
\end{table}

\section{Verilog Implementation}

\begin{lstlisting}[style=verilog, caption={ALU Implementation}, label={lst:alu}]
module alu (
    input wire [63:0] operand_a,
    input wire [63:0] operand_b,
    input wire [3:0] alu_op,
    output reg [63:0] result,
    output reg zero_flag,
    output reg negative_flag,
    output reg carry_flag,
    output reg overflow_flag
);

// ALU operation codes
localparam ALU_ADD = 4'b0000;
localparam ALU_SUB = 4'b0001;
localparam ALU_MUL = 4'b0010;
localparam ALU_DIV = 4'b0011;
localparam ALU_MOD = 4'b0100;
localparam ALU_AND = 4'b0101;
localparam ALU_OR  = 4'b0110;
localparam ALU_XOR = 4'b0111;
localparam ALU_SHL = 4'b1000;
localparam ALU_SHR = 4'b1001;
localparam ALU_SAR = 4'b1010;
localparam ALU_CMP = 4'b1011;

reg [64:0] temp_result; // For carry detection

always @(*) begin
    // Default flag values
    carry_flag = 1'b0;
    overflow_flag = 1'b0;
    
    case (alu_op)
        ALU_ADD: begin
            temp_result = {1'b0, operand_a} + {1'b0, operand_b};
            result = temp_result[63:0];
            carry_flag = temp_result[64];
            overflow_flag = (operand_a[63] == operand_b[63]) && 
                          (result[63] != operand_a[63]);
        end
        
        ALU_SUB: begin
            temp_result = {1'b0, operand_a} - {1'b0, operand_b};
            result = temp_result[63:0];
            carry_flag = temp_result[64];
            overflow_flag = (operand_a[63] != operand_b[63]) && 
                          (result[63] != operand_a[63]);
        end
        
        ALU_MUL: begin
            result = operand_a * operand_b;
        end
        
        ALU_DIV: begin
            result = (operand_b != 0) ? (operand_a / operand_b) : 64'h0;
        end
        
        ALU_MOD: begin
            result = (operand_b != 0) ? (operand_a % operand_b) : 64'h0;
        end
        
        ALU_AND: begin
            result = operand_a & operand_b;
        end
        
        ALU_OR: begin
            result = operand_a | operand_b;
        end
        
        ALU_XOR: begin
            result = operand_a ^ operand_b;
        end
        
        ALU_SHL: begin
            result = operand_a << operand_b[5:0]; // Limit shift to 0-63
        end
        
        ALU_SHR: begin
            result = operand_a >> operand_b[5:0];
        end
        
        ALU_SAR: begin
            result = $signed(operand_a) >>> operand_b[5:0];
        end
        
        ALU_CMP: begin
            temp_result = {1'b0, operand_a} - {1'b0, operand_b};
            result = temp_result[63:0];
        end
        
        default: begin
            result = 64'h0;
        end
    endcase
    
    // Update flags
    zero_flag = (result == 64'h0);
    negative_flag = result[63];
end

endmodule
\end{lstlisting}

\section{Flag Register}

The ALU sets four condition flags after each operation:

\begin{description}
    \item[Zero Flag (ZF):] Set when result is zero
    \item[Negative Flag (NF):] Set when result is negative (bit 63 = 1)
    \item[Carry Flag (CF):] Set on unsigned overflow
    \item[Overflow Flag (OF):] Set on signed overflow
\end{description}

These flags are used by conditional branch instructions to make decisions.

% ============================================================================
% CHAPTER 5: CONTROL UNIT
% ============================================================================
\chapter{Control Unit}

\section{Control Unit Architecture}

The control unit decodes instructions and generates control signals for the datapath. It is implemented as a hardwired control unit rather than microcode for simplicity and performance.

\section{Control Signals}

Table \ref{tab:control-signals} lists all control signals:

\begin{table}[H]
\centering
\caption{Control Signals}
\label{tab:control-signals}
\small
\begin{tabularx}{\textwidth}{|l|l|X|}
\hline
\textbf{Signal} & \textbf{Width} & \textbf{Purpose} \\
\hline
reg\_write & 1 & Enable register write \\
alu\_op & 4 & ALU operation code \\
alu\_src & 1 & ALU source: 0=reg, 1=imm \\
mem\_read & 1 & Enable memory read \\
mem\_write & 1 & Enable memory write \\
mem\_to\_reg & 1 & Write memory to register \\
branch & 1 & Branch instruction \\
jump & 1 & Unconditional jump \\
pc\_src & 2 & PC source: 00=PC+4, 01=branch, 10=jump, 11=register \\
\hline
\end{tabularx}
\end{table}

\section{Verilog Implementation}

\begin{lstlisting}[style=verilog, caption={Control Unit Implementation}, label={lst:control}]
module control_unit (
    input wire [7:0] opcode,
    output reg reg_write,
    output reg [3:0] alu_op,
    output reg alu_src,
    output reg mem_read,
    output reg mem_write,
    output reg mem_to_reg,
    output reg branch,
    output reg jump,
    output reg [1:0] pc_src
);

// Opcode definitions
localparam OP_NOP    = 8'h00;
localparam OP_HALT   = 8'h01;
localparam OP_ADD    = 8'h02;
localparam OP_ADDI   = 8'h03;
localparam OP_SUB    = 8'h04;
localparam OP_SUBI   = 8'h05;
localparam OP_MUL    = 8'h06;
localparam OP_DIV    = 8'h08;
localparam OP_AND    = 8'h0C;
localparam OP_ANDI   = 8'h0D;
localparam OP_OR     = 8'h0E;
localparam OP_XOR    = 8'h10;
localparam OP_SHL    = 8'h12;
localparam OP_SHR    = 8'h14;
localparam OP_LOAD   = 8'h20;
localparam OP_STORE  = 8'h21;
localparam OP_JMP    = 8'h30;
localparam OP_JZ     = 8'h31;
localparam OP_JNZ    = 8'h32;
localparam OP_JGT    = 8'h33;
localparam OP_JLT    = 8'h34;
localparam OP_CALL   = 8'h39;
localparam OP_RET    = 8'h3A;
localparam OP_CMP    = 8'h40;
localparam OP_PUSH   = 8'h42;
localparam OP_POP    = 8'h43;

always @(*) begin
    // Default values
    reg_write = 1'b0;
    alu_op = 4'b0000;
    alu_src = 1'b0;
    mem_read = 1'b0;
    mem_write = 1'b0;
    mem_to_reg = 1'b0;
    branch = 1'b0;
    jump = 1'b0;
    pc_src = 2'b00;
    
    case (opcode)
        OP_NOP: begin
            // No operation
        end
        
        OP_ADD: begin
            reg_write = 1'b1;
            alu_op = 4'b0000; // ADD
            alu_src = 1'b0;   // Register source
        end
        
        OP_ADDI: begin
            reg_write = 1'b1;
            alu_op = 4'b0000; // ADD
            alu_src = 1'b1;   // Immediate source
        end
        
        OP_SUB: begin
            reg_write = 1'b1;
            alu_op = 4'b0001; // SUB
            alu_src = 1'b0;
        end
        
        OP_SUBI: begin
            reg_write = 1'b1;
            alu_op = 4'b0001; // SUB
            alu_src = 1'b1;
        end
        
        OP_MUL: begin
            reg_write = 1'b1;
            alu_op = 4'b0010; // MUL
            alu_src = 1'b0;
        end
        
        OP_DIV: begin
            reg_write = 1'b1;
            alu_op = 4'b0011; // DIV
            alu_src = 1'b0;
        end
        
        OP_AND: begin
            reg_write = 1'b1;
            alu_op = 4'b0101; // AND
            alu_src = 1'b0;
        end
        
        OP_OR: begin
            reg_write = 1'b1;
            alu_op = 4'b0110; // OR
            alu_src = 1'b0;
        end
        
        OP_XOR: begin
            reg_write = 1'b1;
            alu_op = 4'b0111; // XOR
            alu_src = 1'b0;
        end
        
        OP_SHL: begin
            reg_write = 1'b1;
            alu_op = 4'b1000; // SHL
            alu_src = 1'b0;
        end
        
        OP_SHR: begin
            reg_write = 1'b1;
            alu_op = 4'b1001; // SHR
            alu_src = 1'b0;
        end
        
        OP_LOAD: begin
            reg_write = 1'b1;
            mem_read = 1'b1;
            mem_to_reg = 1'b1;
            alu_op = 4'b0000; // ADD for address calculation
            alu_src = 1'b1;   // Use immediate offset
        end
        
        OP_STORE: begin
            mem_write = 1'b1;
            alu_op = 4'b0000; // ADD for address calculation
            alu_src = 1'b1;
        end
        
        OP_JMP: begin
            jump = 1'b1;
            pc_src = 2'b10;
        end
        
        OP_JZ, OP_JNZ, OP_JGT, OP_JLT: begin
            branch = 1'b1;
            pc_src = 2'b01;
        end
        
        OP_CALL: begin
            reg_write = 1'b1; // Write return address to R31
            jump = 1'b1;
            pc_src = 2'b10;
        end
        
        OP_RET: begin
            jump = 1'b1;
            pc_src = 2'b11; // PC from R31
        end
        
        OP_CMP: begin
            alu_op = 4'b1011; // CMP
            alu_src = 1'b0;
        end
        
        OP_PUSH: begin
            mem_write = 1'b1;
            reg_write = 1'b1; // Update SP
        end
        
        OP_POP: begin
            mem_read = 1'b1;
            reg_write = 1'b1;
            mem_to_reg = 1'b1;
        end
        
        default: begin
            // Invalid opcode - all signals remain at default
        end
    endcase
end

endmodule
\end{lstlisting}

% ============================================================================
% CHAPTER 6: MEMORY SUBSYSTEM
% ============================================================================
\chapter{Memory Subsystem}

\section{Memory Architecture}

The A1M architecture uses a Harvard architecture with separate instruction and data memories. This allows simultaneous instruction fetch and data access, improving pipeline performance.

\subsection{Memory Organization}

\begin{itemize}
    \item \textbf{Address Width:} 64 bits (theoretical 16 exabytes)
    \item \textbf{Data Width:} 64 bits (8 bytes per word)
    \item \textbf{Instruction Memory:} Read-only during execution
    \item \textbf{Data Memory:} Read/Write
    \item \textbf{Byte Addressable:} Each byte has a unique address
\end{itemize}

\subsection{Memory Map}

Table \ref{tab:memory-map} shows the default memory layout:

\begin{table}[H]
\centering
\caption{A1M Memory Map}
\label{tab:memory-map}
\begin{tabularx}{\textwidth}{|l|l|X|}
\hline
\textbf{Address Range} & \textbf{Region} & \textbf{Purpose} \\
\hline
0x0000000000000000 - 0x00000000000FFFFF & Code & Instruction memory (1 MB) \\
0x0000000000100000 - 0x00000000001FFFFF & Data & Global/static data (1 MB) \\
0x0000000000200000 - 0x00000000002FFFFF & Heap & Dynamic allocation (1 MB) \\
0x00000000FFFF0000 - 0x00000000FFFFFFFF & Stack & Stack (grows down, 64 KB) \\
\hline
\end{tabularx}
\end{table}

\section{Verilog Implementation}

\begin{lstlisting}[style=verilog, caption={Data Memory Module}, label={lst:data-mem}]
module data_memory #(
    parameter ADDR_WIDTH = 16,  // 64K words = 512KB
    parameter DATA_WIDTH = 64
)(
    input wire clk,
    input wire mem_read,
    input wire mem_write,
    input wire [ADDR_WIDTH-1:0] address,
    input wire [DATA_WIDTH-1:0] write_data,
    output reg [DATA_WIDTH-1:0] read_data
);

// Memory array (64K 64-bit words = 512KB)
reg [DATA_WIDTH-1:0] memory [0:(1 << ADDR_WIDTH)-1];

// Initialize memory to zero
integer i;
initial begin
    for (i = 0; i < (1 << ADDR_WIDTH); i = i + 1) begin
        memory[i] = 64'h0;
    end
end

// Synchronous write, asynchronous read
always @(posedge clk) begin
    if (mem_write) begin
        memory[address] <= write_data;
    end
end

always @(*) begin
    if (mem_read) begin
        read_data = memory[address];
    end else begin
        read_data = 64'h0;
    end
end

endmodule
\end{lstlisting}

\begin{lstlisting}[style=verilog, caption={Instruction Memory Module}, label={lst:inst-mem}]
module instruction_memory #(
    parameter ADDR_WIDTH = 16,
    parameter DATA_WIDTH = 32
)(
    input wire [ADDR_WIDTH-1:0] address,
    output reg [DATA_WIDTH-1:0] instruction
);

// Instruction memory array
reg [DATA_WIDTH-1:0] memory [0:(1 << ADDR_WIDTH)-1];

// Load program from file (simulation)
initial begin
    $readmemh("program.hex", memory);
end

// Asynchronous read
always @(*) begin
    instruction = memory[address];
end

endmodule
\end{lstlisting}

\section{Cache Design (Future Enhancement)}

The current implementation uses simple memory modules. A production design would include:

\begin{itemize}
    \item \textbf{L1 I-Cache:} 32 KB, 4-way set associative
    \item \textbf{L1 D-Cache:} 32 KB, 4-way set associative
    \item \textbf{L2 Unified Cache:} 256 KB, 8-way set associative
    \item \textbf{Cache Line Size:} 64 bytes
    \item \textbf{Write Policy:} Write-back with write-allocate
\end{itemize}

% ============================================================================
% CHAPTER 7: PIPELINE ARCHITECTURE
% ============================================================================
\chapter{Pipeline Architecture}

\section{Five-Stage Pipeline}

The A1M processor implements a classic 5-stage RISC pipeline:

\begin{enumerate}
    \item \textbf{IF (Instruction Fetch):} Fetch instruction from memory
    \item \textbf{ID (Instruction Decode):} Decode instruction and read registers
    \item \textbf{EX (Execute):} Execute ALU operation
    \item \textbf{MEM (Memory):} Access data memory
    \item \textbf{WB (Write Back):} Write result to register
\end{enumerate}

\subsection{Pipeline Diagram}

\begin{verbatim}
Time ->  1    2    3    4    5    6    7
Inst 1: IF   ID   EX  MEM   WB
Inst 2:      IF   ID   EX  MEM   WB
Inst 3:           IF   ID   EX  MEM   WB
Inst 4:                IF   ID   EX  MEM   WB
\end{verbatim}

In ideal conditions, the pipeline achieves a CPI (Cycles Per Instruction) of 1.0 after the initial latency.

\section{Pipeline Hazards}

\subsection{Data Hazards}

Data hazards occur when an instruction depends on the result of a previous instruction still in the pipeline.

\textbf{Example:}
\begin{lstlisting}[style=assembly]
ADD R1, R2, R3    ; R1 = R2 + R3
SUB R4, R1, R5    ; Needs R1 (not yet written!)
\end{lstlisting}

\textbf{Solutions:}
\begin{itemize}
    \item \textbf{Forwarding:} Pass result directly from EX/MEM to EX stage
    \item \textbf{Stalling:} Insert NOP bubbles when forwarding is insufficient
\end{itemize}

\subsection{Control Hazards}

Control hazards occur with branch and jump instructions.

\textbf{Solutions:}
\begin{itemize}
    \item \textbf{Branch Prediction:} Predict branch not taken (simple)
    \item \textbf{Branch Delay Slot:} Execute next instruction regardless
    \item \textbf{Pipeline Flush:} Cancel instructions if prediction wrong
\end{itemize}

\section{Forwarding Unit}

\begin{lstlisting}[style=verilog, caption={Forwarding Unit}, label={lst:forward}]
module forwarding_unit (
    // EX/MEM pipeline register
    input wire [4:0] ex_mem_rd,
    input wire ex_mem_reg_write,
    
    // MEM/WB pipeline register
    input wire [4:0] mem_wb_rd,
    input wire mem_wb_reg_write,
    
    // ID/EX pipeline register
    input wire [4:0] id_ex_rs1,
    input wire [4:0] id_ex_rs2,
    
    // Forwarding control
    output reg [1:0] forward_a,  // 00=no forward, 01=EX/MEM, 10=MEM/WB
    output reg [1:0] forward_b
);

always @(*) begin
    // Default: no forwarding
    forward_a = 2'b00;
    forward_b = 2'b00;
    
    // EX hazard (forward from EX/MEM)
    if (ex_mem_reg_write && (ex_mem_rd != 5'h0)) begin
        if (ex_mem_rd == id_ex_rs1)
            forward_a = 2'b01;
        if (ex_mem_rd == id_ex_rs2)
            forward_b = 2'b01;
    end
    
    // MEM hazard (forward from MEM/WB)
    if (mem_wb_reg_write && (mem_wb_rd != 5'h0)) begin
        if ((mem_wb_rd == id_ex_rs1) && 
            !(ex_mem_reg_write && (ex_mem_rd == id_ex_rs1)))
            forward_a = 2'b10;
        if ((mem_wb_rd == id_ex_rs2) && 
            !(ex_mem_reg_write && (ex_mem_rd == id_ex_rs2)))
            forward_b = 2'b10;
    end
end

endmodule
\end{lstlisting}

\section{Hazard Detection Unit}

\begin{lstlisting}[style=verilog, caption={Hazard Detection Unit}, label={lst:hazard}]
module hazard_detection (
    // ID/EX pipeline register
    input wire id_ex_mem_read,
    input wire [4:0] id_ex_rd,
    
    // IF/ID pipeline register
    input wire [4:0] if_id_rs1,
    input wire [4:0] if_id_rs2,
    
    // Control signals
    output reg stall,
    output reg pc_write,
    output reg if_id_write
);

always @(*) begin
    // Default: no stall
    stall = 1'b0;
    pc_write = 1'b1;
    if_id_write = 1'b1;
    
    // Load-use hazard detection
    if (id_ex_mem_read && 
        ((id_ex_rd == if_id_rs1) || (id_ex_rd == if_id_rs2))) begin
        // Stall the pipeline
        stall = 1'b1;
        pc_write = 1'b0;      // Don't update PC
        if_id_write = 1'b0;   // Don't update IF/ID
    end
end

endmodule
\end{lstlisting}

% ============================================================================
% PART II: ASSEMBLY LANGUAGE
% ============================================================================
\part{Assembly Language Programming}

% ============================================================================
% CHAPTER 8: ASSEMBLY LANGUAGE FUNDAMENTALS
% ============================================================================
\chapter{Assembly Language Fundamentals}

\section{Introduction to A1M Assembly}

Assembly language provides a human-readable representation of machine code. Each assembly instruction corresponds directly to a machine instruction.

\subsection{Assembly Syntax}

The A1M assembly language follows this general syntax:

\begin{verbatim}
label:  INSTRUCTION  operands    ; comment
\end{verbatim}

\textbf{Example:}
\begin{lstlisting}[style=assembly, caption={Simple Assembly Program}]
; Hello World program
.section .data
msg:    .string "Hello, A1M!\n"
len:    .quad 12

.section .text
.global _start

_start:
    ; Write system call
    MOVI R1, 1          ; syscall number (write)
    MOVI R2, 1          ; file descriptor (stdout)
    LOAD R3, msg        ; buffer address
    LOAD R4, len        ; message length
    SYSCALL 0           ; invoke system call
    
    ; Exit system call
    MOVI R1, 0          ; syscall number (exit)
    MOVI R2, 0          ; exit code
    SYSCALL 0
\end{lstlisting}

\section{Directives}

Assembly directives provide instructions to the assembler (not the CPU):

\begin{table}[H]
\centering
\caption{Assembly Directives}
\label{tab:directives}
\begin{tabularx}{\textwidth}{|l|X|}
\hline
\textbf{Directive} & \textbf{Purpose} \\
\hline
.section .text & Begin code section \\
.section .data & Begin data section \\
.section .bss & Begin uninitialized data section \\
.global label & Make label globally visible \\
.byte value & Define byte (8 bits) \\
.word value & Define word (32 bits) \\
.quad value & Define quadword (64 bits) \\
.string "text" & Define null-terminated string \\
.space n & Reserve n bytes of space \\
.align n & Align to n-byte boundary \\
\hline
\end{tabularx}
\end{table}

\section{Labels and Symbols}

Labels mark positions in code or data:

\begin{lstlisting}[style=assembly]
loop_start:
    ADDI R10, R10, 1
    CMPI R10, 100
    JLT loop_start      ; Jump to label
\end{lstlisting}

% ============================================================================
% CHAPTER 9: DATA MOVEMENT INSTRUCTIONS
% ============================================================================
\chapter{Data Movement Instructions}

\section{Register Operations}

\subsection{MOV - Move Register}

Copy value from one register to another:

\begin{lstlisting}[style=assembly]
MOV R1, R2          ; R1 = R2
\end{lstlisting}

Implemented as: \texttt{ADD R1, R2, R0}

\subsection{MOVI - Move Immediate}

Load immediate value into register:

\begin{lstlisting}[style=assembly]
MOVI R1, 42         ; R1 = 42
MOVI R2, -100       ; R2 = -100
\end{lstlisting}

\section{Memory Operations}

\subsection{LOAD - Load from Memory}

Load 64-bit value from memory:

\begin{lstlisting}[style=assembly]
LOAD R1, 0(R3)      ; R1 = Memory[R3 + 0]
LOAD R2, 16(R3)     ; R2 = Memory[R3 + 16]
LOAD R4, array      ; R4 = Memory[array address]
\end{lstlisting}

\subsection{STORE - Store to Memory}

Store 64-bit value to memory:

\begin{lstlisting}[style=assembly]
STORE R1, 0(R3)     ; Memory[R3 + 0] = R1
STORE R2, 16(R3)    ; Memory[R3 + 16] = R2
\end{lstlisting}

\subsection{LOADB/STOREB - Byte Operations}

Load or store single byte:

\begin{lstlisting}[style=assembly]
LOADB R1, 0(R3)     ; R1 = (byte)Memory[R3]
STOREB R2, 0(R4)    ; Memory[R4] = (byte)R2
\end{lstlisting}

\section{Stack Operations}

\subsection{PUSH - Push to Stack}

\begin{lstlisting}[style=assembly]
PUSH R1             ; SP -= 8; Memory[SP] = R1
\end{lstlisting}

Equivalent to:
\begin{lstlisting}[style=assembly]
SUBI R3, R3, 8      ; R3 (SP) -= 8
STORE R1, 0(R3)     ; Memory[SP] = R1
\end{lstlisting}

\subsection{POP - Pop from Stack}

\begin{lstlisting}[style=assembly]
POP R1              ; R1 = Memory[SP]; SP += 8
\end{lstlisting}

\section{Example: Array Access}

\begin{lstlisting}[style=assembly, caption={Array Access Example}]
.section .data
array:  .quad 10, 20, 30, 40, 50

.section .text
_start:
    ; Load base address
    MOVI R10, array     ; R10 = &array[0]
    
    ; Access array[2] (third element)
    MOVI R11, 2         ; Index
    MOVI R12, 8         ; Size of quad (8 bytes)
    MUL R11, R11, R12   ; Offset = index * size
    ADD R13, R10, R11   ; Address = base + offset
    LOAD R14, 0(R13)    ; R14 = array[2] = 30
\end{lstlisting}

% ============================================================================
% CHAPTER 10: ARITHMETIC OPERATIONS
% ============================================================================
\chapter{Arithmetic Operations}

\section{Basic Arithmetic}

\subsection{Addition}

\begin{lstlisting}[style=assembly]
ADD R1, R2, R3      ; R1 = R2 + R3 (register)
ADDI R1, R2, 100    ; R1 = R2 + 100 (immediate)
\end{lstlisting}

\subsection{Subtraction}

\begin{lstlisting}[style=assembly]
SUB R1, R2, R3      ; R1 = R2 - R3
SUBI R1, R2, 50     ; R1 = R2 - 50
\end{lstlisting}

\subsection{Multiplication}

\begin{lstlisting}[style=assembly]
MUL R1, R2, R3      ; R1 = R2 * R3
MULI R1, R2, 10     ; R1 = R2 * 10
\end{lstlisting}

\subsection{Division and Modulo}

\begin{lstlisting}[style=assembly]
DIV R1, R2, R3      ; R1 = R2 / R3 (quotient)
MOD R1, R2, R3      ; R1 = R2 % R3 (remainder)

DIVI R1, R10, 10    ; R1 = R10 / 10
MODI R2, R10, 10    ; R2 = R10 % 10
\end{lstlisting}

\section{Logical Operations}

\subsection{Bitwise AND, OR, XOR}

\begin{lstlisting}[style=assembly]
AND R1, R2, R3      ; R1 = R2 & R3
ANDI R1, R2, 0xFF   ; R1 = R2 & 0xFF (mask lower 8 bits)

OR R1, R2, R3       ; R1 = R2 | R3
ORI R1, R2, 0x80    ; R1 = R2 | 0x80 (set bit 7)

XOR R1, R2, R3      ; R1 = R2 ^ R3
XORI R1, R1, -1     ; R1 = ~R1 (NOT operation)
\end{lstlisting}

\subsection{Shift Operations}

\begin{lstlisting}[style=assembly]
SHL R1, R2, R3      ; R1 = R2 << R3 (left shift)
SHLI R1, R2, 3      ; R1 = R2 << 3

SHR R1, R2, R3      ; R1 = R2 >> R3 (logical right)
SHRI R1, R2, 2      ; R1 = R2 >> 2

SAR R1, R2, R3      ; R1 = R2 >> R3 (arithmetic right)
SARI R1, R2, 4      ; R1 = R2 >> 4 (preserve sign)
\end{lstlisting}

\section{Example: Calculate Expression}

Calculate: $result = (a + b) \times c - d / 2$

\begin{lstlisting}[style=assembly, caption={Complex Expression}]
; Assume: R10=a, R11=b, R12=c, R13=d

; Step 1: temp1 = a + b
ADD R14, R10, R11       ; R14 = a + b

; Step 2: temp2 = temp1 * c
MUL R15, R14, R12       ; R15 = (a + b) * c

; Step 3: temp3 = d / 2
DIVI R16, R13, 2        ; R16 = d / 2

; Step 4: result = temp2 - temp3
SUB R1, R15, R16        ; R1 = (a+b)*c - d/2
\end{lstlisting}

% ============================================================================
% CHAPTER 11: CONTROL FLOW
% ============================================================================
\chapter{Control Flow Instructions}

\section{Unconditional Jumps}

\subsection{JMP - Jump}

Transfer control unconditionally:

\begin{lstlisting}[style=assembly]
JMP target          ; PC = target address
\end{lstlisting}

\subsection{CALL - Function Call}

Call a function (saves return address):

\begin{lstlisting}[style=assembly]
CALL function       ; R31 = PC + 4; PC = function
\end{lstlisting}

\subsection{RET - Return}

Return from function:

\begin{lstlisting}[style=assembly]
RET                 ; PC = R31
\end{lstlisting}

\section{Conditional Branches}

\subsection{Comparison}

Compare two values and set flags:

\begin{lstlisting}[style=assembly]
CMP R1, R2          ; Set flags based on R1 - R2
CMPI R1, 100        ; Set flags based on R1 - 100
\end{lstlisting}

\subsection{Conditional Jumps}

\begin{lstlisting}[style=assembly]
JZ label            ; Jump if Zero Flag set
JNZ label           ; Jump if Zero Flag clear
JGT label           ; Jump if Greater Than (signed)
JLT label           ; Jump if Less Than (signed)
JGE label           ; Jump if Greater or Equal
JLE label           ; Jump if Less or Equal
\end{lstlisting}

\section{Control Flow Examples}

\subsection{If-Then-Else}

\begin{lstlisting}[style=assembly, caption={If-Then-Else Structure}]
; if (R10 > 50) {
;     R1 = 1;
; } else {
;     R1 = 0;
; }

    CMPI R10, 50        ; Compare R10 with 50
    JLE else_block      ; If R10 <= 50, go to else
    
then_block:
    MOVI R1, 1          ; R1 = 1
    JMP end_if          ; Skip else block
    
else_block:
    MOVI R1, 0          ; R1 = 0
    
end_if:
    ; Continue...
\end{lstlisting}

\subsection{While Loop}

\begin{lstlisting}[style=assembly, caption={While Loop}]
; int sum = 0;
; int i = 0;
; while (i < 10) {
;     sum += i;
;     i++;
; }

    MOVI R10, 0         ; sum = 0
    MOVI R11, 0         ; i = 0
    
while_start:
    CMPI R11, 10        ; Compare i with 10
    JGE while_end       ; If i >= 10, exit loop
    
    ADD R10, R10, R11   ; sum += i
    ADDI R11, R11, 1    ; i++
    
    JMP while_start     ; Repeat loop
    
while_end:
    ; R10 contains sum
\end{lstlisting}

\subsection{For Loop}

\begin{lstlisting}[style=assembly, caption={For Loop}]
; for (int i = 0; i < 10; i++) {
;     array[i] = i * 2;
; }

    MOVI R10, 0             ; i = 0
    MOVI R11, array         ; R11 = &array[0]
    
for_start:
    CMPI R10, 10            ; i < 10?
    JGE for_end             ; Exit if i >= 10
    
    ; Calculate i * 2
    MULI R12, R10, 2        ; R12 = i * 2
    
    ; Calculate offset
    MULI R13, R10, 8        ; R13 = i * 8 (byte offset)
    ADD R14, R11, R13       ; R14 = &array[i]
    
    ; Store value
    STORE R12, 0(R14)       ; array[i] = i * 2
    
    ADDI R10, R10, 1        ; i++
    JMP for_start
    
for_end:
\end{lstlisting}

% ============================================================================
% CHAPTER 12: SYSTEM CALLS
% ============================================================================
\chapter{System Calls}

\section{System Call Interface}

System calls provide access to operating system services. Arguments are passed in registers R1-R4.

\subsection{System Call Convention}

\begin{itemize}
    \item \textbf{R1:} System call number
    \item \textbf{R2-R4:} Arguments (up to 3)
    \item \textbf{R1:} Return value
\end{itemize}

\begin{lstlisting}[style=assembly]
SYSCALL 0           ; Invoke system call
\end{lstlisting}

\section{Standard System Calls}

Table \ref{tab:syscalls} lists available system calls:

\begin{table}[H]
\centering
\caption{A1M System Calls}
\label{tab:syscalls}
\small
\begin{tabularx}{\textwidth}{|l|l|X|}
\hline
\textbf{Number} & \textbf{Name} & \textbf{Description} \\
\hline
0 & exit & Exit program (R2 = exit code) \\
1 & write & Write to file (R2=fd, R3=buf, R4=len) \\
2 & read & Read from file (R2=fd, R3=buf, R4=len) \\
3 & open & Open file (R2=path, R3=flags) \\
4 & close & Close file (R2=fd) \\
5 & malloc & Allocate memory (R2=size) → R1=address \\
6 & free & Free memory (R2=address) \\
10 & print\_int & Print integer (R2=value) \\
11 & print\_str & Print string (R2=address) \\
12 & read\_int & Read integer → R1=value \\
13 & read\_str & Read string (R2=buffer, R3=max\_len) \\
\hline
\end{tabularx}
\end{table}

\section{System Call Examples}

\subsection{Print String}

\begin{lstlisting}[style=assembly, caption={Print String Example}]
.section .data
message: .string "Hello, World!\n"

.section .text
    MOVI R1, 11         ; syscall: print_str
    MOVI R2, message    ; address of string
    SYSCALL 0           ; invoke
\end{lstlisting}

\subsection{Read Integer}

\begin{lstlisting}[style=assembly, caption={Read Integer Example}]
    MOVI R1, 12         ; syscall: read_int
    SYSCALL 0           ; invoke
    ; R1 now contains the input value
    MOV R10, R1         ; Save to R10
\end{lstlisting}

\subsection{File I/O}

\begin{lstlisting}[style=assembly, caption={File I/O Example}]
.section .data
filename: .string "output.txt"
buffer:   .string "Hello, File!\n"

.section .text
    ; Open file
    MOVI R1, 3          ; syscall: open
    MOVI R2, filename   ; file path
    MOVI R3, 1          ; flags: O_WRONLY
    SYSCALL 0
    MOV R10, R1         ; Save file descriptor
    
    ; Write to file
    MOVI R1, 1          ; syscall: write
    MOV R2, R10         ; file descriptor
    MOVI R3, buffer     ; data to write
    MOVI R4, 13         ; length
    SYSCALL 0
    
    ; Close file
    MOVI R1, 4          ; syscall: close
    MOV R2, R10         ; file descriptor
    SYSCALL 0
\end{lstlisting}

% ============================================================================
% CHAPTER 13: ASSEMBLER IMPLEMENTATION
% ============================================================================
\chapter{Assembler Implementation}

\section{Two-Pass Assembly}

The A1M assembler uses a two-pass algorithm:

\subsection{Pass 1: Symbol Table Construction}

\begin{enumerate}
    \item Read each line of assembly code
    \item Track current address counter
    \item Record labels and their addresses in symbol table
    \item Calculate sizes of data declarations
\end{enumerate}

\subsection{Pass 2: Code Generation}

\begin{enumerate}
    \item Re-read assembly code
    \item Resolve label references using symbol table
    \item Generate machine code for each instruction
    \item Output binary file
\end{enumerate}

\section{Assembler Architecture}

\begin{lstlisting}[style=c, caption={Assembler Main Structure}]
#include <stdio.h>
#include <stdlib.h>
#include <string.h>

#define MAX_LABELS 1000
#define MAX_LINE 256

typedef struct {
    char name[64];
    uint64_t address;
} Symbol;

Symbol symbol_table[MAX_LABELS];
int symbol_count = 0;

// Function prototypes
void pass1(FILE *input);
void pass2(FILE *input, FILE *output);
uint32_t encode_instruction(char *mnemonic, char *operands);
uint64_t lookup_symbol(char *name);

int main(int argc, char **argv) {
    if (argc != 3) {
        fprintf(stderr, "Usage: %s input.asm output.bin\n", 
                argv[0]);
        return 1;
    }
    
    FILE *input = fopen(argv[1], "r");
    if (!input) {
        perror("Error opening input file");
        return 1;
    }
    
    // Pass 1: Build symbol table
    pass1(input);
    rewind(input);
    
    // Pass 2: Generate code
    FILE *output = fopen(argv[2], "wb");
    if (!output) {
        perror("Error opening output file");
        fclose(input);
        return 1;
    }
    
    pass2(input, output);
    
    fclose(input);
    fclose(output);
    
    printf("Assembly successful: %s -> %s\n", 
           argv[1], argv[2]);
    return 0;
}
\end{lstlisting}

\section{Instruction Encoding}

\begin{lstlisting}[style=c, caption={Instruction Encoder}]
uint32_t encode_r_type(uint8_t opcode, uint8_t rd, 
                       uint8_t rs1, uint8_t rs2) {
    uint32_t instruction = 0;
    instruction |= (opcode & 0xFF) << 24;
    instruction |= (rd & 0x1F) << 19;
    instruction |= (rs1 & 0x1F) << 14;
    instruction |= (rs2 & 0x1F) << 9;
    return instruction;
}

uint32_t encode_i_type(uint8_t opcode, uint8_t rd, 
                       uint8_t rs1, int16_t imm) {
    uint32_t instruction = 0;
    instruction |= (opcode & 0xFF) << 24;
    instruction |= (rd & 0x1F) << 19;
    instruction |= (rs1 & 0x1F) << 14;
    instruction |= (imm & 0x3FFF);
    return instruction;
}

uint32_t encode_j_type(uint8_t opcode, uint32_t address) {
    uint32_t instruction = 0;
    instruction |= (opcode & 0xFF) << 24;
    instruction |= (address & 0xFFFFFF);
    return instruction;
}
\end{lstlisting}

\section{Example: Complete Assembler Usage}

\begin{lstlisting}[style=assembly, caption={Complete Assembly Program}]
; Factorial Calculator
; Calculates n! recursively

.section .data
prompt:  .string "Enter number: "
result:  .string "Factorial: "

.section .text
.global _start

_start:
    ; Print prompt
    MOVI R1, 11
    MOVI R2, prompt
    SYSCALL 0
    
    ; Read number
    MOVI R1, 12
    SYSCALL 0
    MOV R4, R1          ; Save input in R4
    
    ; Calculate factorial
    CALL factorial
    
    ; Print result message
    MOVI R1, 11
    MOVI R2, result
    SYSCALL 0
    
    ; Print result value
    MOVI R1, 10
    MOV R2, R1
    SYSCALL 0
    
    ; Exit
    MOVI R1, 0
    MOVI R2, 0
    SYSCALL 0

; Factorial function
; Input: R4 = n
; Output: R1 = n!
factorial:
    PUSH R31            ; Save return address
    PUSH R4             ; Save n
    
    ; Base case: if (n <= 1) return 1
    CMPI R4, 1
    JGT recursive_case
    
    MOVI R1, 1          ; Return 1
    JMP factorial_end
    
recursive_case:
    ; Recursive case: n * factorial(n-1)
    SUBI R4, R4, 1      ; n - 1
    CALL factorial      ; R1 = factorial(n-1)
    
    POP R4              ; Restore original n
    MUL R1, R1, R4      ; R1 = n * factorial(n-1)
    
factorial_end:
    POP R31             ; Restore return address
    RET
\end{lstlisting}

Assemble with:
\begin{verbatim}
$ ./assembler factorial.asm factorial.bin
\end{verbatim}

% ============================================================================
% PART III: EMULATOR
% ============================================================================
\part{Emulator Implementation}

% ============================================================================
% CHAPTER 14: EMULATOR ARCHITECTURE
% ============================================================================
\chapter{Emulator Architecture}

\section{Emulator Overview}

The A1M emulator is a software implementation of the A1M CPU that executes binary programs. It provides a platform for testing and debugging programs before FPGA synthesis.

\subsection{Features}

\begin{itemize}
    \item Cycle-accurate execution
    \item Full instruction set support
    \item System call emulation
    \item Debugging support (breakpoints, single-step)
    \item Memory inspection
    \item Register dump
\end{itemize}

\section{Emulator Core Structure}

\begin{lstlisting}[style=c, caption={Emulator Core}]
#include <stdio.h>
#include <stdint.h>
#include <stdbool.h>
#include <string.h>

#define MEMORY_SIZE (1024 * 1024)  // 1 MB
#define NUM_REGISTERS 32

typedef struct {
    uint64_t registers[NUM_REGISTERS];
    uint64_t pc;                // Program counter
    uint8_t *memory;
    bool zero_flag;
    bool negative_flag;
    bool carry_flag;
    bool overflow_flag;
    bool halted;
    uint64_t cycle_count;
} CPU;

// Initialize CPU
void cpu_init(CPU *cpu) {
    memset(cpu->registers, 0, sizeof(cpu->registers));
    cpu->pc = 0;
    cpu->memory = calloc(MEMORY_SIZE, 1);
    cpu->zero_flag = false;
    cpu->negative_flag = false;
    cpu->carry_flag = false;
    cpu->overflow_flag = false;
    cpu->halted = false;
    cpu->cycle_count = 0;
}

// Load program into memory
bool load_program(CPU *cpu, const char *filename) {
    FILE *f = fopen(filename, "rb");
    if (!f) {
        perror("Error opening file");
        return false;
    }
    
    size_t bytes_read = fread(cpu->memory, 1, 
                              MEMORY_SIZE, f);
    fclose(f);
    
    printf("Loaded %zu bytes\n", bytes_read);
    return true;
}

// Fetch instruction
uint32_t fetch(CPU *cpu) {
    uint32_t instruction = 0;
    memcpy(&instruction, &cpu->memory[cpu->pc], 4);
    cpu->pc += 4;
    return instruction;
}

// Decode instruction
void decode(uint32_t instruction, uint8_t *opcode,
            uint8_t *rd, uint8_t *rs1, uint8_t *rs2,
            int16_t *imm, uint32_t *addr) {
    *opcode = (instruction >> 24) & 0xFF;
    *rd = (instruction >> 19) & 0x1F;
    *rs1 = (instruction >> 14) & 0x1F;
    *rs2 = (instruction >> 9) & 0x1F;
    *imm = instruction & 0x3FFF;
    // Sign extend immediate
    if (*imm & 0x2000) {
        *imm |= 0xC000;
    }
    *addr = instruction & 0xFFFFFF;
}

// Execute one instruction
void execute(CPU *cpu) {
    if (cpu->halted) return;
    
    uint32_t instruction = fetch(cpu);
    
    uint8_t opcode, rd, rs1, rs2;
    int16_t imm;
    uint32_t addr;
    
    decode(instruction, &opcode, &rd, &rs1, &rs2, 
           &imm, &addr);
    
    // Execute based on opcode
    execute_instruction(cpu, opcode, rd, rs1, rs2, 
                       imm, addr);
    
    // R0 always zero
    cpu->registers[0] = 0;
    
    cpu->cycle_count++;
}
\end{lstlisting}

% ============================================================================
% CHAPTER 15: SYSTEM CALL IMPLEMENTATION
% ============================================================================
\chapter{System Call Implementation}

\section{System Call Handler}

The emulator implements system calls to provide I/O and OS functionality during program execution.

\begin{lstlisting}[style=c, caption={System Call Handler}]
void handle_syscall(CPU *cpu) {
    uint64_t syscall_num = cpu->registers[1];
    uint64_t arg1 = cpu->registers[2];
    uint64_t arg2 = cpu->registers[3];
    uint64_t arg3 = cpu->registers[4];
    
    switch (syscall_num) {
        case 0: // exit
            printf("\nProgram exited with code %llu\n", arg1);
            cpu->halted = true;
            break;
            
        case 1: // write
            syscall_write(cpu, arg1, arg2, arg3);
            break;
            
        case 2: // read
            syscall_read(cpu, arg1, arg2, arg3);
            break;
            
        case 10: // print_int
            printf("%lld", (int64_t)arg1);
            fflush(stdout);
            break;
            
        case 11: // print_str
            syscall_print_str(cpu, arg1);
            break;
            
        case 12: // read_int
            {
                int64_t value;
                scanf("%lld", &value);
                cpu->registers[1] = value;
            }
            break;
            
        default:
            fprintf(stderr, "Unknown syscall: %llu\n", 
                   syscall_num);
    }
}
\end{lstlisting}

% ============================================================================
% PART IV: ALGO HIGH-LEVEL LANGUAGE
% ============================================================================
\part{Algo Programming Language}

% ============================================================================
% CHAPTER 16: INTRODUCTION TO ALGO
% ============================================================================
\chapter{Introduction to Algo Language}

\section{Language Overview}

Algo is a C-like, statically-typed programming language designed for the A1M architecture. It implements a comprehensive subset of C features including modern control flow, operators, and type declarations.

\subsection{Design Goals}

\begin{enumerate}[itemsep=8pt]
    \item \textbf{C Compatibility:} Syntax closely follows C language standards
    \item \textbf{Type Safety:} Static type checking with compile-time constants
    \item \textbf{Efficiency:} Direct compilation to A1M assembly (no runtime overhead)
    \item \textbf{Modern Features:} Includes goto, switch/case, do-while, ternary, ++/--
    \item \textbf{Educational:} Clear mapping between high-level and assembly code
\end{enumerate}

\subsection{Hello World}

\begin{lstlisting}[style=algo, caption={Hello World Program}]
// First Algo program
func main() -> void {
    print("Hello, World!\n");
    print("Welcome to Algo programming!\n");
}
\end{lstlisting}

Compilation and execution:
\begin{lstlisting}[style=bash]
$ ./compiler/algo-compiler hello.algo -o hello.asm
$ ./assembler/a1m-asm hello.asm -o hello.bin
$ ./emulator/a1m-emu hello.bin
Hello, World!
Welcome to Algo programming!
\end{lstlisting}

\section{Lexical Structure}

\subsection{Comments}

\begin{lstlisting}[style=algo]
// Single-line comment (C++ style)

/* Multi-line comment
   Can span multiple lines
   C-style block comment */
\end{lstlisting}

\subsection{Identifiers}

Identifiers must start with a letter or underscore, followed by letters, digits, or underscores:

\begin{lstlisting}[style=algo]
let counter: int = 0;        // Valid
let _temp: int = 5;          // Valid
let value123: int = 999;     // Valid
let my_variable: int = 42;   // Valid
// let 123abc: int = 0;      // Invalid - starts with digit
\end{lstlisting}

\subsection{Keywords}

Reserved keywords in Algo (36 C keywords recognized):

\begin{table}[H]
\centering
\caption{Algo Reserved Keywords}
\small
\begin{tabularx}{\textwidth}{|X|X|X|X|}
\hline
\textbf{Types} & \textbf{Control Flow} & \textbf{Operators} & \textbf{Other} \\
\hline
int, long, char & if, else & sizeof & func, return \\
bool, void & while, do, for & ++, -- & let (variable) \\
struct, union & switch, case & +=, -=, *=, /= & enum, typedef \\
signed, unsigned & default, break & \&=, |=, \^{}=, <<= & const, static \\
short, float & continue, goto & >>=, \%=, ?: & extern, register \\
double & & & auto, volatile \\
\hline
\end{tabularx}
\end{table}

\section{Data Types}

\subsection{Primitive Types}

\begin{table}[H]
\centering
\caption{Algo Data Types (64-bit Architecture)}
\begin{tabularx}{\textwidth}{|l|l|l|X|}
\hline
\textbf{Type} & \textbf{Size} & \textbf{Range} & \textbf{Description} \\
\hline
int & 64 bits & $-2^{63}$ to $2^{63}-1$ & Signed integer (default) \\
long & 64 bits & $-2^{63}$ to $2^{63}-1$ & Alias for int \\
char & 8 bits & 0 to 255 & ASCII character \\
bool & 8 bits & 0 or 1 & Boolean (true/false) \\
void & N/A & N/A & Function return type \\
\hline
\end{tabularx}
\end{table}

\textbf{Note:} All types are 64-bit in registers. \code{sizeof(int) = 8}, \code{sizeof(char) = 1}, \code{sizeof(bool) = 1}.

\subsection{Variable Declarations}

Algo uses the \code{let} keyword for variable declarations:

\begin{lstlisting}[style=algo, caption={Variable Declarations}]
// Basic declarations with let keyword
let x: int;                    // Uninitialized
let age: int = 25;             // Initialized
let grade: char = 'A';         // Character
let is_valid: bool = 1;        // Boolean (1 = true, 0 = false)

// Type inference (type required in current implementation)
let count: int = 0;
let message: int = 42;
\end{lstlisting}

\subsection{Constants and Enumerations}

\begin{lstlisting}[style=algo, caption={Compile-Time Constants with Enums}]
// Enum declarations (compile-time constants)
enum Status {
    IDLE,        // Auto: 0
    RUNNING,     // Auto: 1
    PAUSED,      // Auto: 2
    STOPPED      // Auto: 3
}

enum ErrorCode {
    SUCCESS = 0,
    ERROR = -1,
    TIMEOUT = -2
}

// Using enum constants
let status: int = RUNNING;        // status = 1
let err: int = SUCCESS;           // err = 0

// Enum constants are immutable
// RUNNING = 5;  // ERROR: Cannot assign to constant
\end{lstlisting}

\subsection{Literals}

\begin{lstlisting}[style=algo, caption={Literal Values}]
// Integer literals  
let decimal: int = 42;
let negative: int = -100;

// Character literals
let letter: char = 'A';
let newline: char = '\n';
let tab: char = '\t';

// String literals
print("Hello, World\n");

// Boolean values (0 or 1)
let flag1: int = 1;  // true
let flag2: int = 0;  // false
bool flag2 = false;
\end{lstlisting}

\section{Type Conversion}

\subsection{Implicit Conversion}

\begin{lstlisting}[style=algo]
char c = 65;
int x = c;      // char to int (implicit)
// x == 65

bool b = 1;
int y = b;      // bool to int
// y == 1
\end{lstlisting}

\subsection{Explicit Casting}

\begin{lstlisting}[style=algo]
int x = 1000;
char c = (char)x;  // Truncates to 232 (1000 % 256)

int y = 42;
bool b = (bool)y;  // Non-zero becomes true
\end{lstlisting}

\section{Program Structure}

\subsection{Minimal Program}

Every Algo program must have a \code{main} function:

\begin{lstlisting}[style=algo]
fun main() {
    // Program starts here
}
\end{lstlisting}

\subsection{Complete Program with Functions}

\begin{lstlisting}[style=algo, caption={Complete Algo Program Structure}]
// Global variables
int counter = 0;

// Function declaration
fun increment() {
    counter = counter + 1;
}

fun get_counter() -> int {
    return counter;
}

// Main entry point
fun main() {
    increment();
    increment();
    increment();
    
    int value = get_counter();
    print_int(value);  // Prints: 3
}
\end{lstlisting}

% ============================================================================
% CHAPTER 17: EXPRESSIONS AND OPERATORS
% ============================================================================
\chapter{Expressions and Operators}

\section{Arithmetic Operators}

\subsection{Basic Arithmetic}

\begin{table}[H]
\centering
\caption{Arithmetic Operators (All Implemented)}
\begin{tabularx}{\textwidth}{|l|l|X|l|}
\hline
\textbf{Operator} & \textbf{Name} & \textbf{Description} & \textbf{Example} \\
\hline
+ & Addition & Add two values & \code{a + b} \\
- & Subtraction & Subtract second from first & \code{a - b} \\
* & Multiplication & Multiply two values & \code{a * b} \\
/ & Division & Integer division & \code{a / b} \\
\% & Modulo & Remainder after division & \code{a \% b} \\
++ & Increment & Add 1 (prefix/postfix) & \code{++a}, \code{a++} \\
-- & Decrement & Subtract 1 (prefix/postfix) & \code{--a}, \code{a--} \\
\hline
\end{tabularx}
\end{table}

\begin{lstlisting}[style=algo, caption={Arithmetic Examples}]
func arithmetic_demo() -> void {
    let a: int = 10;
    let b: int = 3;
    
    let sum: int = a + b;         // 13
    let diff: int = a - b;        // 7
    let product: int = a * b;     // 30
    let quotient: int = a / b;    // 3 (integer division)
    let remainder: int = a % b;   // 1
    
    // Complex expressions
    let result: int = (a + b) * 2 - a / b;
    // result = (10 + 3) * 2 - 10 / 3
    // result = 13 * 2 - 3 = 23
}
\end{lstlisting}

\subsection{Increment and Decrement Operators}

\begin{lstlisting}[style=algo, caption={Prefix vs Postfix ++/--}]
func increment_demo() -> void {
    let x: int = 5;
    
    // Prefix: increment THEN use value
    let a: int = ++x;  // x=6, a=6
    
    // Postfix: use value THEN increment  
    let b: int = x++;  // b=6, x=7
    
    // Decrement works similarly
    let c: int = --x;  // x=6, c=6
    let d: int = x--;  // d=6, x=5
    
    // Common in loops
    let i: int = 0;
    while (i < 10) {
        print(i);
        i++;  // Equivalent to: i = i + 1
    }
}
\end{lstlisting}

\subsection{Compound Assignment Operators}

\begin{table}[H]
\centering
\caption{Compound Assignment Operators (All 10 Implemented)}
\begin{tabularx}{\textwidth}{|l|X|l|}
\hline
\textbf{Operator} & \textbf{Equivalent To} & \textbf{Example} \\
\hline
+= & \code{x = x + y} & \code{x += 5} \\
-= & \code{x = x - y} & \code{x -= 3} \\
*= & \code{x = x * y} & \code{x *= 2} \\
/= & \code{x = x / y} & \code{x /= 4} \\
\%= & \code{x = x \% y} & \code{x \%= 10} \\
\&= & \code{x = x \& y} & \code{x \&= 0xFF} \\
|= & \code{x = x | y} & \code{x |= 0x80} \\
\^{}= & \code{x = x \^{} y} & \code{x \^{}= 1} \\
<<= & \code{x = x << y} & \code{x <<= 2} \\
>>= & \code{x = x >> y} & \code{x >>= 1} \\
\hline
\end{tabularx}
\end{table}

\begin{lstlisting}[style=algo, caption={Compound Assignment Examples}]
func compound_demo() -> void {
    let x: int = 10;
    
    x += 5;   // x = 15
    x -= 7;   // x = 8
    x *= 4;   // x = 32
    x /= 2;   // x = 16
    x %= 5;   // x = 1
    
    let flags: int = 0;
    flags |= 8;    // Set bit 3
    flags &= 15;   // Mask to low 4 bits
    flags ^= 7;    // Toggle bits 0-2
    
    let shift: int = 4;
    shift <<= 2;   // shift = 16 (multiply by 4)
    shift >>= 3;   // shift = 2 (divide by 8)
}
\end{lstlisting}

\subsection{Operator Precedence}

\begin{table}[H]
\centering
\caption{Operator Precedence (Highest to Lowest)}
\small
\begin{tabular}{|l|l|l|}
\hline
\textbf{Precedence} & \textbf{Operators} & \textbf{Associativity} \\
\hline
1 (Highest) & \code{()}, function call, \code{sizeof} & Left \\
2 & \code{!}, \code{~}, \code{++}, \code{--}, unary \code{-} & Right \\
3 & \code{*}, \code{/}, \code{\%} & Left \\
4 & \code{+}, \code{-} & Left \\
5 & \code{<<}, \code{>>} & Left \\
6 & \code{<}, \code{<=}, \code{>}, \code{>=} & Left \\
7 & \code{==}, \code{!=} & Left \\
8 & \code{\&} (bitwise AND) & Left \\
9 & \code{\^{}} (bitwise XOR) & Left \\
10 & \code{|} (bitwise OR) & Left \\
11 & \code{\&\&} (logical AND) & Left \\
12 & \code{||} (logical OR) & Left \\
13 & \code{? :} (ternary) & Right \\
14 (Lowest) & \code{=}, \code{+=}, \code{-=}, etc. & Right \\
\hline
\end{tabular}
\end{table}

\section{Comparison Operators}

\begin{table}[H]
\centering
\caption{Comparison Operators (All Implemented)}
\begin{tabularx}{\textwidth}{|l|l|X|}
\hline
\textbf{Operator} & \textbf{Name} & \textbf{Returns 1 (true) if} \\
\hline
== & Equal & Operands are equal \\
!= & Not equal & Operands are not equal \\
< & Less than & Left < Right \\
> & Greater than & Left > Right \\
<= & Less or equal & Left $\leq$ Right \\
>= & Greater or equal & Left $\geq$ Right \\
\hline
\end{tabularx}
\end{table}

\begin{lstlisting}[style=algo, caption={Comparison Examples}]
func comparison_demo() -> void {
    let x: int = 10;
    let y: int = 20;
    
    let eq: int = (x == y);     // 0 (false)
    let ne: int = (x != y);     // 1 (true)
    let lt: int = (x < y);      // 1 (true)
    let gt: int = (x > y);      // 0 (false)
    let le: int = (x <= y);     // 1 (true)
    let ge: int = (x >= y);     // 0 (false)
    
    // Used in conditions
    if (x < y) {
        print("x is less than y\n");
    }
}
\end{lstlisting}

\section{Logical Operators}

\begin{table}[H]
\centering
\caption{Logical Operators (Short-Circuit Evaluation)}
\begin{tabularx}{\textwidth}{|l|l|X|}
\hline
\textbf{Operator} & \textbf{Name} & \textbf{Description} \\
\hline
\&\& & Logical AND & True if both operands are non-zero \\
|| & Logical OR & True if at least one operand is non-zero \\
! & Logical NOT & Returns 0 if operand is non-zero, 1 if zero \\
\hline
\end{tabularx}
\end{table}

\begin{lstlisting}[style=algo, caption={Logical Operations with Short-Circuit}]
func logical_demo() -> void {
    let a: int = 1;  // true
    let b: int = 0;  // false
    
    let and_result: int = a && b;   // 0 (false)
    let or_result: int = a || b;    // 1 (true)
    let not_a: int = !a;            // 0 (false)
    let not_b: int = !b;            // 1 (true)
    
    // Complex logical expressions
    let x: int = 10;
    let y: int = 20;
    let z: int = 15;
    
    // Check if z is between x and y
    let in_range: int = (z > x) && (z < y);  // 1 (true)
    
    // Short-circuit: second not evaluated if first is false
    if (x != 0 && y / x > 2) {
        print("Safe division\n");
    }
}
\end{lstlisting}

\section{Bitwise Operators}

\begin{table}[H]
\centering
\caption{Bitwise Operators (All Implemented)}
\begin{tabularx}{\textwidth}{|l|l|X|}
\hline
\textbf{Operator} & \textbf{Name} & \textbf{Description} \\
\hline
\& & Bitwise AND & Performs AND on each bit \\
| & Bitwise OR & Performs OR on each bit \\
\^{} & Bitwise XOR & Performs exclusive OR on each bit \\
\~{} & Bitwise NOT & Inverts all bits \\
<< & Left shift & Shifts bits left (multiply by 2\^{}n) \\
>> & Right shift & Shifts bits right (divide by 2\^{}n) \\
\hline
\end{tabularx}
\end{table}

\begin{lstlisting}[style=algo, caption={Bitwise Operations}]
func bitwise_demo() -> void {
    let x: int = 42;  // Binary: 0b00101010
    
    print("x = ");
    print(x);
    print("\n");
    
    print("x << 1 = ");
    print(x << 1);  // 84 (multiply by 2)
    print("\n");
    
    print("x >> 1 = ");
    print(x >> 1);  // 21 (divide by 2)
    print("\n");
    
    print("x & 15 = ");
    print(x & 15);  // 10 (extract low 4 bits)
    print("\n");
    
    print("x | 128 = ");
    print(x | 128);  // 170 (set bit 7)
    print("\n");
    
    print("~x = ");
    print(~x);  // -43 (bitwise complement)
    print("\n");
}
\end{lstlisting}

\section{Ternary Conditional Operator}

\begin{lstlisting}[style=algo, caption={Ternary Operator (? :)}]
func max(a: int, b: int) -> int {
    return (a > b) ? a : b;
}

func abs(x: int) -> int {
    return (x < 0) ? -x : x;
}

func classify(score: int) -> int {
    // Nested ternary (chain)
    return (score >= 90) ? 1 : 
           (score >= 70) ? 2 : 
           (score >= 50) ? 3 : 4;
}

func ternary_demo() -> void {
    let result: int = (10 > 5) ? 100 : 200;  // result = 100
    print("Max(15, 23) = ");
    print(max(15, 23));  // 23
    print("\n");
}
\end{lstlisting}

\section{sizeof Operator}

\begin{lstlisting}[style=algo, caption={sizeof - Compile-Time Operator}]
func sizeof_demo() -> void {
    print("sizeof(int) = ");
    print(sizeof(int));    // 8 (64-bit)
    print("\n");
    
    print("sizeof(char) = ");
    print(sizeof(char));   // 1
    print("\n");
    
    print("sizeof(bool) = ");
    print(sizeof(bool));   // 1
    print("\n");
    
    let x: int = 42;
    print("sizeof(x) = ");
    print(sizeof(x));      // 8
    print("\n");
    
    print("sizeof(10 + 20) = ");
    print(sizeof(10 + 20)); // 8 (expression)
    print("\n");
}
\end{lstlisting}

\section{Expression Examples}

\subsection{Complex Mathematical Expressions}

\begin{lstlisting}[style=algo, caption={Mathematical Formulas}]
// Quadratic formula: x = (-b +/- sqrt(b^2-4ac)) / 2a
func quadratic_positive(a: int, b: int, c: int) -> int {
    let discriminant: int = b * b - 4 * a * c;
    let result: int = (-b + discriminant) / (2 * a);
    return result;
}

// Area of circle: A = pi*r^2
func circle_area(radius: int) -> int {
    let pi_approx: int = 314;  // pi ~ 3.14 * 100
    let area: int = (pi_approx * radius * radius) / 100;
    return area;
}

// Factorial: n! = n * (n-1) * ... * 1
func factorial(n: int) -> int {
    if (n <= 1) {
        return 1;
    }
    return n * factorial(n - 1);
}
\end{lstlisting}

% ============================================================================
% CHAPTER 18: CONTROL STRUCTURES
% ============================================================================
\chapter{Control Structures}

\section{If Statement}

\subsection{Basic If}

\begin{lstlisting}[style=algo, caption={Simple If Statement}]
func temperature_check() -> void {
    let temperature: int = 25;
    
    if (temperature > 30) {
        print("It's hot!\n");
    }
    
    // Execution continues here
}
\end{lstlisting}

\subsection{If-Else}

\begin{lstlisting}[style=algo, caption={If-Else Statement}]
func check_even(n: int) -> void {
    if (n % 2 == 0) {
        print("Even number\n");
    } else {
        print("Odd number\n");
    }
}

func main() -> void {
    check_even(42);  // Prints: Even number
    check_even(17);  // Prints: Odd number
}
\end{lstlisting}

\subsection{If-Else-If Chains}

\begin{lstlisting}[style=algo, caption={Multiple Conditions}]
func grade_to_letter(score: int) -> void {
    if (score >= 90) {
        print("Grade: A\n");
    } else if (score >= 80) {
        print("Grade: B\n");
    } else if (score >= 70) {
        print("Grade: C\n");
    } else if (score >= 60) {
        print("Grade: D\n");
    } else {
        print("Grade: F\n");
    }
}

func main() -> void {
    grade_to_letter(85);  // Prints: Grade: B
    grade_to_letter(55);  // Prints: Grade: F
}
\end{lstlisting}

\subsection{Nested If Statements}

\begin{lstlisting}[style=algo, caption={Nested Conditions}]
func check_number(x: int) -> void {
    if (x > 0) {
        if (x % 2 == 0) {
            print("Positive even\n");
        } else {
            print("Positive odd\n");
        }
    } else if (x < 0) {
        if (x % 2 == 0) {
            print("Negative even\n");
        } else {
            print("Negative odd\n");
        }
    } else {
        print("Zero\n");
    }
}
\end{lstlisting}

\section{While Loop}

\subsection{Basic While Loop}

\begin{lstlisting}[style=algo, caption={Counting with While}]
func count_to_ten() -> void {
    let i: int = 1;
    while (i <= 10) {
        print(i);
        print("\n");
        i++;
    }
}
\end{lstlisting}

\subsection{While with Complex Conditions}

\begin{lstlisting}[style=algo, caption={Fibonacci Numbers}]
func fibonacci(max: int) -> void {
    let a: int = 0;
    let b: int = 1;
    
    print("Fibonacci series:\n");
    while (a <= max) {
        print(a);
        print(" ");
        
        let temp: int = a + b;
        a = b;
        b = temp;
    }
    print("\n");
}
\end{lstlisting}

\section{Do-While Loop}

The do-while loop executes at least once, then checks the condition.

\begin{lstlisting}[style=algo, caption={Do-While Loop}]
func menu_system() -> void {
    let choice: int = 0;
    
    do {
        print("Menu:\n");
        print("1. Option A\n");
        print("2. Option B\n");
        print("3. Exit\n");
        
        choice = get_input();  // Assume this function exists
        
        if (choice == 1) {
            print("Selected A\n");
        } else if (choice == 2) {
            print("Selected B\n");
        }
    } while (choice != 3);
    
    print("Goodbye!\n");
}

func countdown() -> void {
    let i: int = 5;
    
    do {
        print(i);
        print("...\n");
        i--;
    } while (i > 0);
    
    print("Blastoff!\n");
}
\end{lstlisting}

\section{For Loop}

\subsection{Basic For Loop}

\begin{lstlisting}[style=algo, caption={Standard For Loop}]
func print_squares() -> void {
    let i: int = 0;
    for (i = 1; i <= 10; i++) {
        let square: int = i * i;
        print(i);
        print(" squared = ");
        print(square);
        print("\n");
    }
}
\end{lstlisting}

\subsection{For Loop Variations}

\begin{lstlisting}[style=algo, caption={For Loop Examples}]
// Counting backwards
func countdown_for() -> void {
    let i: int = 0;
    for (i = 10; i >= 1; i--) {
        print(i);
        print(" ");
    }
    print("\n");
}

// Step by 2
func even_numbers() -> void {
    let i: int = 0;
    for (i = 0; i <= 20; i += 2) {
        print(i);
        print(" ");
    }
    print("\n");
}

// Multiple variables
func multiple_vars() -> void {
    let i: int = 0;
    let j: int = 0;
    for (i = 0; i < 5; i++) {
        for (j = 0; j < 5; j++) {
            print(i * 5 + j);
            print(" ");
        }
        print("\n");
    }
}
\end{lstlisting}

\section{Switch-Case Statement}

The switch statement provides multi-way branching.

\begin{lstlisting}[style=algo, caption={Switch-Case Example}]
func day_name(day: int) -> void {
    switch (day) {
        case 1:
            print("Monday\n");
            break;
        case 2:
            print("Tuesday\n");
            break;
        case 3:
            print("Wednesday\n");
            break;
        case 4:
            print("Thursday\n");
            break;
        case 5:
            print("Friday\n");
            break;
        case 6:
            print("Saturday\n");
            break;
        case 7:
            print("Sunday\n");
            break;
        default:
            print("Invalid day\n");
            break;
    }
}
\end{lstlisting}

\subsection{Switch with Fall-Through}

\begin{lstlisting}[style=algo, caption={Switch Fall-Through}]
func categorize_grade(grade: int) -> void {
    switch (grade) {
        case 10:
        case 9:
            print("Excellent (A)\n");
            break;
        case 8:
            print("Good (B)\n");
            break;
        case 7:
        case 6:
            print("Average (C)\n");
            break;
        case 5:
            print("Pass (D)\n");
            break;
        default:
            print("Fail (F)\n");
            break;
    }
}

func calculator(op: int, a: int, b: int) -> int {
    let result: int = 0;
    
    switch (op) {
        case 1:  // Addition
            result = a + b;
            break;
        case 2:  // Subtraction
            result = a - b;
            break;
        case 3:  // Multiplication
            result = a * b;
            break;
        case 4:  // Division
            if (b != 0) {
                result = a / b;
            } else {
                print("Error: Division by zero\n");
                result = 0;
            }
            break;
        default:
            print("Invalid operator\n");
            result = 0;
            break;
    }
    
    return result;
}
\end{lstlisting}

\section{Break and Continue}

\subsection{Break Statement}

\begin{lstlisting}[style=algo, caption={Using Break}]
func find_first_divisor(n: int) -> int {
    let i: int = 0;
    for (i = 2; i < n; i++) {
        if (n % i == 0) {
            print("First divisor: ");
            print(i);
            print("\n");
            break;  // Exit loop early
        }
    }
    return i;
}

func search_array(arr: int, size: int, target: int) -> int {
    let i: int = 0;
    let found: int = 0;
    
    for (i = 0; i < size; i++) {
        if (arr == target) {
            found = 1;
            break;
        }
    }
    
    return found;
}
\end{lstlisting}

\subsection{Continue Statement}

\begin{lstlisting}[style=algo, caption={Using Continue}]
func print_odd_numbers() -> void {
    let i: int = 0;
    for (i = 1; i <= 20; i++) {
        if (i % 2 == 0) {
            continue;  // Skip even numbers
        }
        print(i);
        print(" ");
    }
    print("\n");
}

func skip_multiples_of_three() -> void {
    let i: int = 0;
    for (i = 1; i <= 30; i++) {
        if (i % 3 == 0) {
            continue;  // Skip multiples of 3
        }
        print(i);
        print(" ");
    }
    print("\n");
}
\end{lstlisting}

\section{Goto Statement and Labels}

Goto allows unconditional jumps to labels.

\begin{lstlisting}[style=algo, caption={Goto and Labels}]
func goto_example() -> void {
    let i: int = 0;
    
    start:
    if (i >= 5) {
        goto end;
    }
    
    print(i);
    print(" ");
    i++;
    goto start;
    
    end:
    print("\nDone!\n");
}

func error_handling() -> void {
    let status: int = do_operation();
    if (status != 0) {
        goto error;
    }
    
    status = another_operation();
    if (status != 0) {
        goto error;
    }
    
    print("Success!\n");
    goto cleanup;
    
    error:
    print("Error occurred!\n");
    
    cleanup:
    print("Cleaning up...\n");
}
\end{lstlisting}

\subsection{Nested Loop Break with Goto}

\begin{lstlisting}[style=algo, caption={Breaking Nested Loops}]
func find_in_matrix() -> void {
    let found: int = 0;
    let i: int = 0;
    let j: int = 0;
    
    for (i = 0; i < 10; i++) {
        for (j = 0; j < 10; j++) {
            if (i * 10 + j == 42) {
                found = 1;
                goto found_it;
            }
        }
    }
    
    print("Not found\n");
    goto end;
    
    found_it:
    print("Found at: ");
    print(i);
    print(", ");
    print(j);
    print("\n");
    
    end:
    print("Search complete\n");
}
\end{lstlisting}

\section{Comprehensive Control Flow Example}

\begin{lstlisting}[style=algo, caption={Complete Control Flow Demo}]
func control_flow_demo() -> void {
    let choice: int = 0;
    let running: int = 1;
    
    print("Control Flow Demonstration\n");
    
    main_loop:
    while (running) {
        print("\n=== Menu ===\n");
        print("1. If-Else Demo\n");
        print("2. Loop Demo\n");
        print("3. Switch Demo\n");
        print("4. Goto Demo\n");
        print("5. Exit\n");
        print("Choice: ");
        
        choice = 3;  // Simulated input
        
        switch (choice) {
            case 1:
                let x: int = 42;
                if (x > 0) {
                    print("Positive\n");
                } else if (x < 0) {
                    print("Negative\n");
                } else {
                    print("Zero\n");
                }
                break;
                
            case 2:
                let i: int = 0;
                do {
                    print(i);
                    print(" ");
                    i++;
                } while (i < 5);
                print("\n");
                break;
                
            case 3:
                day_name(3);
                break;
                
            case 4:
                goto jump_target;
                
            case 5:
                running = 0;
                break;
                
            default:
                print("Invalid choice\n");
                break;
        }
        
        if (choice == 4) {
            jump_target:
            print("Jumped to target!\n");
        }
        
        // Simulate exit after one iteration
        running = 0;
    }
    
    print("Program terminated.\n");
}
\end{lstlisting}

% ============================================================================
% CHAPTER 19: FUNCTIONS
% ============================================================================
\chapter{Functions and Procedures}

\section{Function Basics}

\subsection{Function Declaration Syntax}

In Algo, functions are declared with the \code{func} keyword (not \code{fun}):

\begin{lstlisting}[style=algo, caption={Function Declaration Syntax}]
// Basic syntax:
func function_name(param: type, ...) -> return_type {
    // function body
}

// No parameters, no return value
func greet() -> void {
    print("Hello, World!\n");
}

// With parameters, no return
func print_sum(a: int, b: int) -> void {
    let sum: int = a + b;
    print(sum);
    print("\n");
}

// With return value
func add(a: int, b: int) -> int {
    return a + b;
}

// Multiple parameters and return
func max(a: int, b: int, c: int) -> int {
    let result: int = a;
    
    if (b > result) {
        result = b;
    }
    
    if (c > result) {
        result = c;
    }
    
    return result;
}
\end{lstlisting}

\subsection{Function Calls}

\begin{lstlisting}[style=algo, caption={Calling Functions}]
func main() -> void {
    // Call function with no return
    greet();
    
    // Call function with parameters
    print_sum(10, 20);
    
    // Use return value
    let sum: int = add(5, 7);
    print(sum);
    print("\n");
    
    // Use in expression
    let result: int = add(3, 4) * 2;
    
    // Nested calls
    let maximum: int = max(add(1, 2), add(3, 4), add(5, 6));
    print("Maximum: ");
    print(maximum);
    print("\n");
}
\end{lstlisting}

\section{Recursion}

\subsection{Factorial}

\begin{lstlisting}[style=algo, caption={Recursive Factorial}]
func factorial(n: int) -> int {
    // Base case
    if (n <= 1) {
        return 1;
    }
    
    // Recursive case
    return n * factorial(n - 1);
}

func main() -> void {
    print("5! = ");
    print(factorial(5));  // 120
    print("\n");
    
    print("10! = ");
    print(factorial(10));  // 3628800
    print("\n");
}
\end{lstlisting}

\subsection{Fibonacci Sequence}

\begin{lstlisting}[style=algo, caption={Recursive Fibonacci}]
func fibonacci(n: int) -> int {
    if (n <= 1) {
        return n;
    }
    return fibonacci(n - 1) + fibonacci(n - 2);
}

func print_fibonacci(count: int) -> void {
    print("Fibonacci sequence:\n");
    let i: int = 0;
    for (i = 0; i < count; i++) {
        print(fibonacci(i));
        print(" ");
    }
    print("\n");
}

func main() -> void {
    print_fibonacci(10);  // 0 1 1 2 3 5 8 13 21 34
}
\end{lstlisting}

\subsection{Greatest Common Divisor (GCD)}

\begin{lstlisting}[style=algo, caption={Euclidean Algorithm}]
func gcd(a: int, b: int) -> int {
    if (b == 0) {
        return a;
    }
    return gcd(b, a % b);
}

func main() -> void {
    print("GCD(48, 18) = ");
    print(gcd(48, 18));  // 6
    print("\n");
    
    print("GCD(100, 35) = ");
    print(gcd(100, 35));  // 5
    print("\n");
}
\end{lstlisting}

\section{Advanced Function Examples}

\subsection{Prime Number Checker}

\begin{lstlisting}[style=algo, caption={Prime Number Functions}]
func is_prime(n: int) -> int {
    if (n <= 1) {
        return 0;  // false
    }
    
    if (n == 2) {
        return 1;  // true
    }
    
    if (n % 2 == 0) {
        return 0;  // false
    }
    
    let i: int = 3;
    while (i * i <= n) {
        if (n % i == 0) {
            return 0;  // false
        }
        i += 2;
    }
    
    return 1;  // true
}

func print_primes(limit: int) -> void {
    print("Prime numbers up to ");
    print(limit);
    print(":\n");
    
    let n: int = 2;
    for (n = 2; n <= limit; n++) {
        if (is_prime(n)) {
            print(n);
            print(" ");
        }
    }
    print("\n");
}

func main() -> void {
    print_primes(50);
}
\end{lstlisting}

\subsection{Power Function}

\begin{lstlisting}[style=algo, caption={Iterative and Recursive Power}]
func power_iterative(base: int, exp: int) -> int {
    let result: int = 1;
    let i: int = 0;
    
    for (i = 0; i < exp; i++) {
        result *= base;
    }
    
    return result;
}

func power_recursive(base: int, exp: int) -> int {
    if (exp == 0) {
        return 1;
    }
    
    if (exp == 1) {
        return base;
    }
    
    // Fast exponentiation
    if (exp % 2 == 0) {
        let half: int = power_recursive(base, exp / 2);
        return half * half;
    } else {
        return base * power_recursive(base, exp - 1);
    }
}

func main() -> void {
    print("2^10 = ");
    print(power_iterative(2, 10));  // 1024
    print("\n");
    
    print("3^5 = ");
    print(power_recursive(3, 5));  // 243
    print("\n");
}
\end{lstlisting}

\section{Function Composition and Higher-Order Patterns}

\begin{lstlisting}[style=algo, caption={Function Composition}]
func square(x: int) -> int {
    return x * x;
}

func double_value(x: int) -> int {
    return x * 2;
}

func apply_twice(x: int, is_square: int) -> int {
    if (is_square) {
        return square(square(x));
    } else {
        return double_value(double_value(x));
    }
}

func main() -> void {
    let x: int = 3;
    
    print("square(3) = ");
    print(square(x));  // 9
    print("\n");
    
    print("square(square(3)) = ");
    print(apply_twice(x, 1));  // 81
    print("\n");
    
    print("double(double(3)) = ");
    print(apply_twice(x, 0));  // 12
    print("\n");
}
\end{lstlisting}

% ============================================================================
% CHAPTER 20: ARRAYS
% ============================================================================
\chapter{Arrays and Data Structures}

\section{Array Basics}

\textbf{Note:} Algo currently does NOT support true heap-allocated arrays. All variables are register-based (64-bit values). This chapter describes conceptual array operations that can be simulated.

\subsection{Simulated Array Operations}

\begin{lstlisting}[style=algo, caption={Array-Style Variables}]
func array_simulation() -> void {
    // Individual variables act like array elements
    let arr0: int = 10;
    let arr1: int = 20;
    let arr2: int = 30;
    let arr3: int = 40;
    let arr4: int = 50;
    
    // Process "array" elements
    print("Elements: ");
    print(arr0); print(" ");
    print(arr1); print(" ");
    print(arr2); print(" ");
    print(arr3); print(" ");
    print(arr4); print("\n");
    
    // Calculate sum
    let sum: int = arr0 + arr1 + arr2 + arr3 + arr4;
    print("Sum: ");
    print(sum);  // 150
    print("\n");
}
\end{lstlisting}

\subsection{Array-Like Processing with Loops}

\begin{lstlisting}[style=algo, caption={Sequential Data Processing}]
func process_sequence(count: int) -> void {
    let i: int = 0;
    let sum: int = 0;
    
    print("Sequence: ");
    for (i = 0; i < count; i++) {
        let value: int = i * i;  // Generate square numbers
        print(value);
        print(" ");
        sum += value;
    }
    print("\n");
    
    print("Sum of squares: ");
    print(sum);
    print("\n");
}

func main() -> void {
    process_sequence(10);  // 0 1 4 9 16 25 36 49 64 81
}
\end{lstlisting}

\section{Searching and Sorting Concepts}

\subsection{Linear Search Simulation}

\begin{lstlisting}[style=algo, caption={Search Among Multiple Variables}]
func linear_search_demo() -> void {
    // Simulated array
    let v0: int = 42;
    let v1: int = 17;
    let v2: int = 93;
    let v3: int = 28;
    let v4: int = 55;
    
    let target: int = 93;
    let found: int = 0;
    let position: int = -1;
    
    // Manual search
    if (v0 == target) {
        found = 1;
        position = 0;
        goto search_done;
    }
    
    if (v1 == target) {
        found = 1;
        position = 1;
        goto search_done;
    }
    
    if (v2 == target) {
        found = 1;
        position = 2;
        goto search_done;
    }
    
    if (v3 == target) {
        found = 1;
        position = 3;
        goto search_done;
    }
    
    if (v4 == target) {
        found = 1;
        position = 4;
    }
    
    search_done:
    if (found) {
        print("Found ");
        print(target);
        print(" at position ");
        print(position);
        print("\n");
    } else {
        print("Not found\n");
    }
}
\end{lstlisting}

\subsection{Sorting Algorithms (Conceptual)}

\begin{lstlisting}[style=algo, caption={Bubble Sort Concept}]
func bubble_sort_demo() -> void {
    // Manual bubble sort with 5 elements
    let a: int = 64;
    let b: int = 34;
    let c: int = 25;
    let d: int = 12;
    let e: int = 22;
    
    print("Before: ");
    print(a); print(" ");
    print(b); print(" ");
    print(c); print(" ");
    print(d); print(" ");
    print(e); print("\n");
    
    // Pass 1
    if (a > b) { let t: int = a; a = b; b = t; }
    if (b > c) { let t: int = b; b = c; c = t; }
    if (c > d) { let t: int = c; c = d; d = t; }
    if (d > e) { let t: int = d; d = e; e = t; }
    
    // Pass 2
    if (a > b) { let t: int = a; a = b; b = t; }
    if (b > c) { let t: int = b; b = c; c = t; }
    if (c > d) { let t: int = c; c = d; d = t; }
    
    // Pass 3
    if (a > b) { let t: int = a; a = b; b = t; }
    if (b > c) { let t: int = b; b = c; c = t; }
    
    // Pass 4
    if (a > b) { let t: int = a; a = b; b = t; }
    
    print("After:  ");
    print(a); print(" ");
    print(b); print(" ");
    print(c); print(" ");
    print(d); print(" ");
    print(e); print("\n");
}
\end{lstlisting}

\section{Matrix Operations (Simulated)}

\begin{lstlisting}[style=algo, caption={2x2 Matrix Multiplication}]
func matrix_multiply_2x2() -> void {
    // Matrix A
    let a00: int = 1; let a01: int = 2;
    let a10: int = 3; let a11: int = 4;
    
    // Matrix B
    let b00: int = 5; let b01: int = 6;
    let b10: int = 7; let b11: int = 8;
    
    // Result C = A * B
    let c00: int = a00 * b00 + a01 * b10;  // 19
    let c01: int = a00 * b01 + a01 * b11;  // 22
    let c10: int = a10 * b00 + a11 * b10;  // 43
    let c11: int = a10 * b01 + a11 * b11;  // 50
    
    print("Result matrix:\n");
    print(c00); print(" "); print(c01); print("\n");
    print(c10); print(" "); print(c11); print("\n");
}
\end{lstlisting}

% ============================================================================
% CHAPTER 21: ENUMERATIONS
% ============================================================================
\chapter{Enumerations and Named Constants}

\section{Enum Declaration}

Enums in Algo define compile-time integer constants:

\begin{lstlisting}[style=algo, caption={Enum Syntax}]
enum Color {
    RED,      // 0
    GREEN,    // 1
    BLUE,     // 2
    YELLOW    // 3
}

enum Status {
    SUCCESS = 0,
    ERROR = -1,
    PENDING = 1,
    TIMEOUT = 2
}

enum FileMode {
    READ = 1,
    WRITE = 2,
    APPEND = 4,
    BINARY = 8
}
\end{lstlisting}

\section{Using Enums}

\begin{lstlisting}[style=algo, caption={Enum Usage Examples}]
enum Direction {
    NORTH,
    SOUTH,
    EAST,
    WEST
}

func move(dir: int) -> void {
    switch (dir) {
        case 0:  // NORTH
            print("Moving north\n");
            break;
        case 1:  // SOUTH
            print("Moving south\n");
            break;
        case 2:  // EAST
            print("Moving east\n");
            break;
        case 3:  // WEST
            print("Moving west\n");
            break;
        default:
            print("Invalid direction\n");
            break;
    }
}

func main() -> void {
    let current_dir: int = 0;  // NORTH
    move(current_dir);
    
    current_dir = 2;  // EAST
    move(current_dir);
}
\end{lstlisting}

\subsection{Enum Flags (Bit Flags)}

\begin{lstlisting}[style=algo, caption={Using Enums as Bit Flags}]
enum Permission {
    NONE = 0,
    READ = 1,
    WRITE = 2,
    EXECUTE = 4,
    DELETE = 8
}

func check_permission(perms: int, required: int) -> int {
    return (perms & required) == required;
}

func main() -> void {
    let user_perms: int = 3;  // READ | WRITE
    
    if (check_permission(user_perms, 1)) {  // READ
        print("Can read\n");
    }
    
    if (check_permission(user_perms, 2)) {  // WRITE
        print("Can write\n");
    }
    
    if (!check_permission(user_perms, 4)) {  // EXECUTE
        print("Cannot execute\n");
    }
    
    // Add execute permission
    user_perms |= 4;
    
    if (check_permission(user_perms, 4)) {
        print("Now can execute\n");
    }
}
\end{lstlisting}

\section{State Machines with Enums}

\begin{lstlisting}[style=algo, caption={State Machine Example}]
enum State {
    IDLE,
    RUNNING,
    PAUSED,
    STOPPED,
    ERROR
}

enum Event {
    START,
    PAUSE,
    RESUME,
    STOP,
    FAIL
}

func process_event(current_state: int, event: int) -> int {
    let next_state: int = current_state;
    
    switch (current_state) {
        case 0:  // IDLE
            if (event == 0) {  // START
                next_state = 1;  // RUNNING
                print("IDLE -> RUNNING\n");
            }
            break;
            
        case 1:  // RUNNING
            if (event == 1) {  // PAUSE
                next_state = 2;  // PAUSED
                print("RUNNING -> PAUSED\n");
            } else if (event == 3) {  // STOP
                next_state = 3;  // STOPPED
                print("RUNNING -> STOPPED\n");
            } else if (event == 4) {  // FAIL
                next_state = 4;  // ERROR
                print("RUNNING -> ERROR\n");
            }
            break;
            
        case 2:  // PAUSED
            if (event == 2) {  // RESUME
                next_state = 1;  // RUNNING
                print("PAUSED -> RUNNING\n");
            } else if (event == 3) {  // STOP
                next_state = 3;  // STOPPED
                print("PAUSED -> STOPPED\n");
            }
            break;
            
        case 3:  // STOPPED
        case 4:  // ERROR
            print("Terminal state\n");
            break;
            
        default:
            print("Unknown state\n");
            break;
    }
    
    return next_state;
}

func main() -> void {
    let state: int = 0;  // IDLE
    
    state = process_event(state, 0);  // START
    state = process_event(state, 1);  // PAUSE
    state = process_event(state, 2);  // RESUME
    state = process_event(state, 3);  // STOP
}
\end{lstlisting}

% ============================================================================
% CHAPTER 22: COMPLETE PROGRAMS
% ============================================================================
\chapter{Complete Algo Programs}

\section{Calculator Program}

\begin{lstlisting}[style=algo, caption={Simple Calculator}]
func calculator() -> void {
    let num1: int = 15;
    let num2: int = 3;
    let operation: int = 1;  // 1=add, 2=sub, 3=mul, 4=div
    let result: int = 0;
    
    print("Calculator: ");
    print(num1);
    
    switch (operation) {
        case 1:
            print(" + ");
            result = num1 + num2;
            break;
        case 2:
            print(" - ");
            result = num1 - num2;
            break;
        case 3:
            print(" * ");
            result = num1 * num2;
            break;
        case 4:
            print(" / ");
            if (num2 != 0) {
                result = num1 / num2;
            } else {
                print("ERROR: Division by zero\n");
                return;
            }
            break;
        default:
            print("ERROR: Invalid operation\n");
            return;
    }
    
    print(num2);
    print(" = ");
    print(result);
    print("\n");
}

func main() -> void {
    calculator();
}
\end{lstlisting}

\section{Number Guessing Game}

\begin{lstlisting}[style=algo, caption={Guess the Number}]
func guess_number_game() -> void {
    let secret: int = 42;
    let guess: int = 0;
    let attempts: int = 0;
    let max_attempts: int = 5;
    let found: int = 0;
    
    print("Guess the number (1-100)\n");
    print("You have ");
    print(max_attempts);
    print(" attempts.\n\n");
    
    // Simulated guesses
    let guesses0: int = 50;
    let guesses1: int = 30;
    let guesses2: int = 40;
    let guesses3: int = 45;
    let guesses4: int = 42;
    
    // Attempt 1
    guess = guesses0;
    attempts++;
    print("Attempt ");
    print(attempts);
    print(": ");
    print(guess);
    
    if (guess == secret) {
        found = 1;
        goto game_over;
    } else if (guess < secret) {
        print(" - Too low!\n");
    } else {
        print(" - Too high!\n");
    }
    
    // Attempt 2
    guess = guesses1;
    attempts++;
    print("Attempt ");
    print(attempts);
    print(": ");
    print(guess);
    
    if (guess == secret) {
        found = 1;
        goto game_over;
    } else if (guess < secret) {
        print(" - Too low!\n");
    } else {
        print(" - Too high!\n");
    }
    
    // Attempt 3
    guess = guesses2;
    attempts++;
    print("Attempt ");
    print(attempts);
    print(": ");
    print(guess);
    
    if (guess == secret) {
        found = 1;
        goto game_over;
    } else if (guess < secret) {
        print(" - Too low!\n");
    } else {
        print(" - Too high!\n");
    }
    
    // Attempt 4
    guess = guesses3;
    attempts++;
    print("Attempt ");
    print(attempts);
    print(": ");
    print(guess);
    
    if (guess == secret) {
        found = 1;
        goto game_over;
    } else if (guess < secret) {
        print(" - Too low!\n");
    } else {
        print(" - Too high!\n");
    }
    
    // Attempt 5
    guess = guesses4;
    attempts++;
    print("Attempt ");
    print(attempts);
    print(": ");
    print(guess);
    
    if (guess == secret) {
        found = 1;
    } else if (guess < secret) {
        print(" - Too low!\n");
    } else {
        print(" - Too high!\n");
    }
    
    game_over:
    print("\n");
    
    if (found) {
        print("Congratulations! You found it in ");
        print(attempts);
        print(" attempts!\n");
    } else {
        print("Game over! The number was ");
        print(secret);
        print("\n");
    }
}

func main() -> void {
    guess_number_game();
}
\end{lstlisting}

\section{Prime Factorization}

\begin{lstlisting}[style=algo, caption={Prime Factors}]
func print_prime_factors(n: int) -> void {
    print("Prime factors of ");
    print(n);
    print(": ");
    
    // Handle factor 2
    while (n % 2 == 0) {
        print("2 ");
        n /= 2;
    }
    
    // Check odd factors
    let i: int = 3;
    while (i * i <= n) {
        while (n % i == 0) {
            print(i);
            print(" ");
            n /= i;
        }
        i += 2;
    }
    
    // If n is still > 1, it's prime
    if (n > 1) {
        print(n);
        print(" ");
    }
    
    print("\n");
}

func main() -> void {
    print_prime_factors(60);    // 2 2 3 5
    print_prime_factors(315);   // 3 3 5 7
    print_prime_factors(128);   // 2 2 2 2 2 2 2
}
\end{lstlisting}

\section{Tower of Hanoi}

\begin{lstlisting}[style=algo, caption={Tower of Hanoi Puzzle}]
func hanoi(n: int, from: int, to: int, aux: int) -> void {
    if (n == 1) {
        print("Move disk 1 from rod ");
        print(from);
        print(" to rod ");
        print(to);
        print("\n");
        return;
    }
    
    hanoi(n - 1, from, aux, to);
    
    print("Move disk ");
    print(n);
    print(" from rod ");
    print(from);
    print(" to rod ");
    print(to);
    print("\n");
    
    hanoi(n - 1, aux, to, from);
}

func main() -> void {
    let num_disks: int = 3;
    print("Tower of Hanoi with ");
    print(num_disks);
    print(" disks:\n\n");
    
    hanoi(num_disks, 1, 3, 2);
}
\end{lstlisting}

\section{Collatz Conjecture}

\begin{lstlisting}[style=algo, caption={Collatz Sequence}]
func collatz(n: int) -> int {
    let steps: int = 0;
    
    print("Collatz sequence for ");
    print(n);
    print(":\n");
    print(n);
    print(" ");
    
    while (n != 1) {
        if (n % 2 == 0) {
            n /= 2;
        } else {
            n = 3 * n + 1;
        }
        
        print(n);
        print(" ");
        steps++;
        
        // Safety limit
        if (steps > 100) {
            print("...\n");
            break;
        }
    }
    
    print("\nSteps: ");
    print(steps);
    print("\n\n");
    
    return steps;
}

func main() -> void {
    collatz(12);
    collatz(27);
    collatz(19);
}
\end{lstlisting}

% ============================================================================
% CHAPTER 23: ADVANCED ALGORITHMS
% ============================================================================
\chapter{Advanced Algorithms}

\section{Binary Search}

\begin{lstlisting}[style=algo, caption={Binary Search (Simulated)}]
func binary_search_demo() -> void {
    // Sorted "array" of 10 elements
    let a0: int = 5;   let a1: int = 12;
    let a2: int = 17;  let a3: int = 23;
    let a4: int = 28;  let a5: int = 34;
    let a6: int = 41;  let a7: int = 55;
    let a8: int = 67;  let a9: int = 89;
    
    let target: int = 34;
    let left: int = 0;
    let right: int = 9;
    let found: int = 0;
    let position: int = -1;
    
    print("Searching for ");
    print(target);
    print("...\n");
    
    while (left <= right) {
        let mid: int = (left + right) / 2;
        let mid_value: int = 0;
        
        // Get value at mid (manual indexing)
        if (mid == 0) mid_value = a0;
        else if (mid == 1) mid_value = a1;
        else if (mid == 2) mid_value = a2;
        else if (mid == 3) mid_value = a3;
        else if (mid == 4) mid_value = a4;
        else if (mid == 5) mid_value = a5;
        else if (mid == 6) mid_value = a6;
        else if (mid == 7) mid_value = a7;
        else if (mid == 8) mid_value = a8;
        else if (mid == 9) mid_value = a9;
        
        print("  Checking index ");
        print(mid);
        print(" (value=");
        print(mid_value);
        print(")\n");
        
        if (mid_value == target) {
            found = 1;
            position = mid;
            break;
        } else if (mid_value < target) {
            left = mid + 1;
        } else {
            right = mid - 1;
        }
    }
    
    if (found) {
        print("Found at index ");
        print(position);
        print("\n");
    } else {
        print("Not found\n");
    }
}

func main() -> void {
    binary_search_demo();
}
\end{lstlisting}

\section{Sieve of Eratosthenes}

\begin{lstlisting}[style=algo, caption={Find All Primes Up To N}]
func sieve_of_eratosthenes(limit: int) -> void {
    // Simulated boolean array (1=prime, 0=composite)
    let p2: int = 1; let p3: int = 1; let p4: int = 1;
    let p5: int = 1; let p6: int = 1; let p7: int = 1;
    let p8: int = 1; let p9: int = 1; let p10: int = 1;
    let p11: int = 1; let p12: int = 1; let p13: int = 1;
    let p14: int = 1; let p15: int = 1; let p16: int = 1;
    let p17: int = 1; let p18: int = 1; let p19: int = 1;
    let p20: int = 1;
    
    print("Primes up to 20:\n");
    
    // Mark multiples of 2
    p4 = 0; p6 = 0; p8 = 0; p10 = 0;
    p12 = 0; p14 = 0; p16 = 0; p18 = 0; p20 = 0;
    
    // Mark multiples of 3
    p6 = 0; p9 = 0; p12 = 0; p15 = 0; p18 = 0;
    
    // Mark multiples of 5
    p10 = 0; p15 = 0; p20 = 0;
    
    // Print primes
    if (p2) print("2 ");
    if (p3) print("3 ");
    if (p5) print("5 ");
    if (p7) print("7 ");
    if (p11) print("11 ");
    if (p13) print("13 ");
    if (p17) print("17 ");
    if (p19) print("19 ");
    print("\n");
}

func main() -> void {
    sieve_of_eratosthenes(20);
}
\end{lstlisting}

\section{Dynamic Programming: Fibonacci}

\begin{lstlisting}[style=algo, caption={Fibonacci with Memoization Concept}]
func fibonacci_iterative(n: int) -> int {
    if (n <= 1) {
        return n;
    }
    
    let prev2: int = 0;
    let prev1: int = 1;
    let current: int = 0;
    let i: int = 2;
    
    for (i = 2; i <= n; i++) {
        current = prev1 + prev2;
        prev2 = prev1;
        prev1 = current;
    }
    
    return current;
}

func compare_fibonacci() -> void {
    print("Fibonacci numbers (iterative):\n");
    
    let i: int = 0;
    for (i = 0; i <= 15; i++) {
        print("F(");
        print(i);
        print(") = ");
        print(fibonacci_iterative(i));
        print("\n");
    }
}

func main() -> void {
    compare_fibonacci();
}
\end{lstlisting}

\section{Greatest Common Divisor Variants}

\begin{lstlisting}[style=algo, caption={GCD and LCM}]
func gcd_iterative(a: int, b: int) -> int {
    while (b != 0) {
        let temp: int = b;
        b = a % b;
        a = temp;
    }
    return a;
}

func lcm(a: int, b: int) -> int {
    return (a * b) / gcd_iterative(a, b);
}

func gcd_multiple(a: int, b: int, c: int) -> int {
    let gcd_ab: int = gcd_iterative(a, b);
    return gcd_iterative(gcd_ab, c);
}

func main() -> void {
    print("GCD(48, 18) = ");
    print(gcd_iterative(48, 18));
    print("\n");
    
    print("LCM(12, 18) = ");
    print(lcm(12, 18));
    print("\n");
    
    print("GCD(24, 36, 60) = ");
    print(gcd_multiple(24, 36, 60));
    print("\n");
}
\end{lstlisting}

% ============================================================================
% CHAPTER 24: BITWISE ALGORITHMS
% ============================================================================
\chapter{Bitwise Algorithms and Tricks}

\section{Bit Manipulation Fundamentals}

\begin{lstlisting}[style=algo, caption={Bit Operations}]
func set_bit(num: int, pos: int) -> int {
    return num | (1 << pos);
}

func clear_bit(num: int, pos: int) -> int {
    return num & ~(1 << pos);
}

func toggle_bit(num: int, pos: int) -> int {
    return num ^ (1 << pos);
}

func check_bit(num: int, pos: int) -> int {
    return (num & (1 << pos)) != 0;
}

func bit_manipulation_demo() -> void {
    let x: int = 42;  // Binary: 00101010
    
    print("Original: ");
    print(x);
    print("\n");
    
    x = set_bit(x, 0);  // Set bit 0
    print("After set bit 0: ");
    print(x);  // 43
    print("\n");
    
    x = clear_bit(x, 1);  // Clear bit 1
    print("After clear bit 1: ");
    print(x);  // 41
    print("\n");
    
    x = toggle_bit(x, 3);  // Toggle bit 3
    print("After toggle bit 3: ");
    print(x);  // 49
    print("\n");
}

func main() -> void {
    bit_manipulation_demo();
}
\end{lstlisting}

\section{Count Set Bits}

\begin{lstlisting}[style=algo, caption={Population Count}]
func count_set_bits(n: int) -> int {
    let count: int = 0;
    
    while (n != 0) {
        count += n & 1;
        n >>= 1;
    }
    
    return count;
}

func count_set_bits_fast(n: int) -> int {
    let count: int = 0;
    
    // Brian Kernighan's algorithm
    while (n != 0) {
        n &= (n - 1);  // Clear rightmost set bit
        count++;
    }
    
    return count;
}

func main() -> void {
    let num: int = 42;  // Binary: 101010
    
    print("Number: ");
    print(num);
    print("\n");
    
    print("Set bits: ");
    print(count_set_bits(num));  // 3
    print("\n");
    
    print("Set bits (fast): ");
    print(count_set_bits_fast(num));  // 3
    print("\n");
}
\end{lstlisting}

\section{Power of Two Check}

\begin{lstlisting}[style=algo, caption={Check if Power of 2}]
func is_power_of_two(n: int) -> int {
    if (n <= 0) {
        return 0;
    }
    return (n & (n - 1)) == 0;
}

func next_power_of_two(n: int) -> int {
    if (n <= 0) {
        return 1;
    }
    
    n--;
    n |= n >> 1;
    n |= n >> 2;
    n |= n >> 4;
    n |= n >> 8;
    n |= n >> 16;
    n |= n >> 32;
    n++;
    
    return n;
}

func main() -> void {
    let tests0: int = 16;
    let tests1: int = 17;
    let tests2: int = 32;
    let tests3: int = 33;
    
    print("Is 16 power of 2? ");
    print(is_power_of_two(tests0));
    print("\n");
    
    print("Is 17 power of 2? ");
    print(is_power_of_two(tests1));
    print("\n");
    
    print("Next power of 2 after 17: ");
    print(next_power_of_two(tests1));
    print("\n");
}
\end{lstlisting}

\section{Swap Without Temporary Variable}

\begin{lstlisting}[style=algo, caption={XOR Swap}]
func xor_swap_demo() -> void {
    let a: int = 42;
    let b: int = 17;
    
    print("Before: a=");
    print(a);
    print(", b=");
    print(b);
    print("\n");
    
    // XOR swap
    a ^= b;
    b ^= a;
    a ^= b;
    
    print("After:  a=");
    print(a);
    print(", b=");
    print(b);
    print("\n");
}

func main() -> void {
    xor_swap_demo();
}
\end{lstlisting}

\section{Find Missing Number}

\begin{lstlisting}[style=algo, caption={Find Missing Number Using XOR}]
func find_missing() -> void {
    // Array contains 0-9 except one number
    let arr0: int = 0;
    let arr1: int = 1;
    let arr2: int = 2;
    let arr3: int = 3;
    let arr4: int = 4;
    // Missing: 5
    let arr5: int = 6;
    let arr6: int = 7;
    let arr7: int = 8;
    let arr8: int = 9;
    
    let xor_all: int = 0;
    let xor_arr: int = 0;
    
    // XOR all numbers 0-9
    let i: int = 0;
    for (i = 0; i <= 9; i++) {
        xor_all ^= i;
    }
    
    // XOR array elements
    xor_arr ^= arr0;
    xor_arr ^= arr1;
    xor_arr ^= arr2;
    xor_arr ^= arr3;
    xor_arr ^= arr4;
    xor_arr ^= arr5;
    xor_arr ^= arr6;
    xor_arr ^= arr7;
    xor_arr ^= arr8;
    
    let missing: int = xor_all ^ xor_arr;
    
    print("Missing number: ");
    print(missing);  // 5
    print("\n");
}

func main() -> void {
    find_missing();
}
\end{lstlisting}

% ============================================================================
% CHAPTER 25: OPTIMIZATION TECHNIQUES
% ============================================================================
\chapter{Optimization and Best Practices}

\section{Loop Optimization}

\begin{lstlisting}[style=algo, caption={Loop Unrolling}]
func sum_array_normal(count: int) -> int {
    let sum: int = 0;
    let i: int = 0;
    
    for (i = 0; i < count; i++) {
        sum += i;
    }
    
    return sum;
}

func sum_array_unrolled(count: int) -> int {
    let sum: int = 0;
    let i: int = 0;
    
    // Process 4 elements at a time
    for (i = 0; i < count; i += 4) {
        sum += i;
        sum += i + 1;
        sum += i + 2;
        sum += i + 3;
    }
    
    return sum;
}

func main() -> void {
    print("Sum (normal): ");
    print(sum_array_normal(100));
    print("\n");
    
    print("Sum (unrolled): ");
    print(sum_array_unrolled(100));
    print("\n");
}
\end{lstlisting}

\section{Strength Reduction}

\begin{lstlisting}[style=algo, caption={Replace Expensive Operations}]
func multiply_by_constant_slow(x: int) -> int {
    return x * 16;  // Multiplication
}

func multiply_by_constant_fast(x: int) -> int {
    return x << 4;  // Left shift (faster)
}

func divide_by_constant_slow(x: int) -> int {
    return x / 8;  // Division
}

func divide_by_constant_fast(x: int) -> int {
    return x >> 3;  // Right shift (faster)
}

func main() -> void {
    let x: int = 100;
    
    print("100 * 16 = ");
    print(multiply_by_constant_fast(x));
    print("\n");
    
    print("100 / 8 = ");
    print(divide_by_constant_fast(x));
    print("\n");
}
\end{lstlisting}

\section{Short-Circuit Evaluation}

\begin{lstlisting}[style=algo, caption={Efficient Conditionals}]
func expensive_check() -> int {
    // Simulate expensive operation
    let i: int = 0;
    let sum: int = 0;
    for (i = 0; i < 1000; i++) {
        sum += i;
    }
    return sum > 0;
}

func cheap_check(x: int) -> int {
    return x > 0;
}

func optimized_condition(x: int) -> void {
    // Put cheap check first
    if (cheap_check(x) && expensive_check()) {
        print("Both conditions true\n");
    }
    
    // If x <= 0, expensive_check() never runs!
}

func main() -> void {
    optimized_condition(5);   // Both checks run
    optimized_condition(-1);  // Only cheap check runs
}
\end{lstlisting}

\section{Table Lookup}

\begin{lstlisting}[style=algo, caption={Precomputed Values}]
func square_computed(x: int) -> int {
    return x * x;
}

func square_lookup(x: int) -> int {
    // Lookup table for squares 0-10
    if (x == 0) return 0;
    if (x == 1) return 1;
    if (x == 2) return 4;
    if (x == 3) return 9;
    if (x == 4) return 16;
    if (x == 5) return 25;
    if (x == 6) return 36;
    if (x == 7) return 49;
    if (x == 8) return 64;
    if (x == 9) return 81;
    if (x == 10) return 100;
    
    // Fall back to computation
    return x * x;
}

func main() -> void {
    let i: int = 0;
    
    print("Squares using lookup:\n");
    for (i = 0; i <= 10; i++) {
        print(i);
        print("^2 = ");
        print(square_lookup(i));
        print("\n");
    }
}
\end{lstlisting}

\section{Early Exit Patterns}

\begin{lstlisting}[style=algo, caption={Avoid Unnecessary Work}]
func find_divisor_slow(n: int) -> int {
    let divisor: int = 0;
    let i: int = 2;
    
    for (i = 2; i < n; i++) {
        if (n % i == 0) {
            divisor = i;
        }
    }
    
    return divisor;  // Always checks all numbers
}

func find_divisor_fast(n: int) -> int {
    let i: int = 2;
    
    for (i = 2; i < n; i++) {
        if (n % i == 0) {
            return i;  // Early exit!
        }
    }
    
    return n;
}

func main() -> void {
    print("First divisor of 1000: ");
    print(find_divisor_fast(1000));  // Returns 2 immediately
    print("\n");
}
\end{lstlisting}

\section{Tail Recursion}

\begin{lstlisting}[style=algo, caption={Tail-Recursive Factorial}]
func factorial_recursive(n: int) -> int {
    if (n <= 1) return 1;
    return n * factorial_recursive(n - 1);  // Not tail-recursive
}

func factorial_tail_helper(n: int, acc: int) -> int {
    if (n <= 1) return acc;
    return factorial_tail_helper(n - 1, n * acc);  // Tail-recursive
}

func factorial_tail(n: int) -> int {
    return factorial_tail_helper(n, 1);
}

func main() -> void {
    print("Factorial(5) = ");
    print(factorial_tail(5));
    print("\n");
    
    print("Factorial(10) = ");
    print(factorial_tail(10));
    print("\n");
}
\end{lstlisting}

\section{Code Size vs Speed Tradeoffs}

\begin{lstlisting}[style=algo, caption={Optimization Strategy}]
// Small code, slower
func abs_branching(x: int) -> int {
    if (x < 0) {
        return -x;
    }
    return x;
}

// Larger code, potentially faster (no branch)
func abs_bitwise(x: int) -> int {
    let mask: int = x >> 63;  // All 1s if negative, all 0s if positive
    return (x + mask) ^ mask;
}

// Ternary operator (compact and clear)
func abs_ternary(x: int) -> int {
    return (x < 0) ? -x : x;
}

func main() -> void {
    let test: int = -42;
    
    print("abs(-42) branching: ");
    print(abs_branching(test));
    print("\n");
    
    print("abs(-42) bitwise:   ");
    print(abs_bitwise(test));
    print("\n");
    
    print("abs(-42) ternary:   ");
    print(abs_ternary(test));
    print("\n");
}
\end{lstlisting}

% CHAPTER 26: STANDARD LIBRARY
% ============================================================================
\chapter{Standard Library and System Functions}

\section{Overview}

The Algo language provides a minimal standard library focused on I/O operations. Most utility functions must be implemented by the user, encouraging understanding of fundamental algorithms.

\section{I/O Functions}

\subsection{Console Output}

\begin{lstlisting}[style=algo, caption={Standard Output Functions}]
// Built-in print function - accepts int values
func print(value: int) -> void;

// Print with newline
func println(value: int) -> void {
    print(value);
    print(10);  // ASCII newline
}

// Example usage
func main() -> void {
    print(42);
    print(10);  // newline
    
    println(100);  // Prints 100 with newline
}
\end{lstlisting}

\subsection{Formatted Output Examples}

\begin{lstlisting}[style=algo, caption={Custom Output Functions}]
// Print multiple values on one line
func print_values(a: int, b: int, c: int) -> void {
    print(a);
    print(32);  // space
    print(b);
    print(32);  // space
    print(c);
    print(10);  // newline
}

// Print with label
func print_labeled(label: int, value: int) -> void {
    print(label);
    print(58);   // ':'
    print(32);   // space
    print(value);
    print(10);   // newline
}

func demo() -> void {
    print_values(10, 20, 30);     // Output: 10 20 30
    print_labeled(65, 42);         // Output: A: 42
}
\end{lstlisting}

\section{String Operations}

\textbf{Note:} Algo currently does not have built-in string support. Strings are simulated using integer arrays where each element represents a character's ASCII value.

\subsection{Character and String Simulation}

\begin{lstlisting}[style=algo, caption={Simulating String Operations}]
// Print a simulated "string" as ASCII codes
func print_chars(codes: int, count: int) -> void {
    let i: int = 0;
    for (i = 0; i < count; i++) {
        print(codes);  // Would print ASCII code
    }
}

// Demonstrate character codes
func char_demo() -> void {
    print(72);   // 'H'
    print(101);  // 'e'
    print(108);  // 'l'
    print(108);  // 'l'
    print(111);  // 'o'
    print(10);   // newline
}
\end{lstlisting}

\section{Mathematical Utilities}

Users must implement mathematical functions as needed.

\begin{lstlisting}[style=algo, caption={User-Implemented Math Functions}]
// Absolute value
func abs(x: int) -> int {
    return (x < 0) ? -x : x;
}

// Maximum of two values
func max(a: int, b: int) -> int {
    return (a > b) ? a : b;
}

// Minimum of two values
func min(a: int, b: int) -> int {
    return (a < b) ? a : b;
}

// Power function (integer exponentiation)
func pow(base: int, exp: int) -> int {
    if (exp == 0) {
        return 1;
    }
    
    let result: int = 1;
    let i: int = 0;
    
    for (i = 0; i < exp; i++) {
        result *= base;
    }
    
    return result;
}

// Integer square root (Newton's method)
func isqrt(n: int) -> int {
    if (n < 2) {
        return n;
    }
    
    let x: int = n;
    let y: int = (x + 1) / 2;
    
    while (y < x) {
        x = y;
        y = (x + n / x) / 2;
    }
    
    return x;
}

// GCD using Euclidean algorithm
func gcd(a: int, b: int) -> int {
    while (b != 0) {
        let temp: int = b;
        b = a % b;
        a = temp;
    }
    return a;
}

// LCM using GCD
func lcm(a: int, b: int) -> int {
    return (a * b) / gcd(a, b);
}
\end{lstlisting}

\section{Array Utilities}

\begin{lstlisting}[style=algo, caption={Common Array Operations}]
// Find maximum in array (register-based simulation)
func array_max(a: int, b: int, c: int, d: int, e: int) -> int {
    let max_val: int = a;
    
    if (b > max_val) { max_val = b; }
    if (c > max_val) { max_val = c; }
    if (d > max_val) { max_val = d; }
    if (e > max_val) { max_val = e; }
    
    return max_val;
}

// Sum of values
func sum5(a: int, b: int, c: int, d: int, e: int) -> int {
    return a + b + c + d + e;
}

// Average of values
func average5(a: int, b: int, c: int, d: int, e: int) -> int {
    return (a + b + c + d + e) / 5;
}

// Linear search simulation
func find_value(val1: int, val2: int, val3: int, 
                val4: int, val5: int, target: int) -> int {
    if (val1 == target) { return 0; }
    if (val2 == target) { return 1; }
    if (val3 == target) { return 2; }
    if (val4 == target) { return 3; }
    if (val5 == target) { return 4; }
    return -1;  // Not found
}
\end{lstlisting}

\section{Bit Manipulation Utilities}

\begin{lstlisting}[style=algo, caption={Bit Operations}]
// Set bit at position
func set_bit(value: int, pos: int) -> int {
    return value | (1 << pos);
}

// Clear bit at position
func clear_bit(value: int, pos: int) -> int {
    return value & ~(1 << pos);
}

// Toggle bit at position
func toggle_bit(value: int, pos: int) -> int {
    return value ^ (1 << pos);
}

// Test if bit is set
func test_bit(value: int, pos: int) -> int {
    return (value & (1 << pos)) ? 1 : 0;
}

// Count set bits (population count)
func popcount(x: int) -> int {
    let count: int = 0;
    
    while (x != 0) {
        count++;
        x &= (x - 1);  // Clear lowest set bit
    }
    
    return count;
}

// Reverse bits in 8-bit value
func reverse_bits8(x: int) -> int {
    let result: int = 0;
    let i: int = 0;
    
    for (i = 0; i < 8; i++) {
        result <<= 1;
        result |= (x & 1);
        x >>= 1;
    }
    
    return result;
}
\end{lstlisting}

\section{Timing and Performance}

\begin{lstlisting}[style=algo, caption={Simple Performance Measurement}]
// Busy-wait delay (CPU cycles)
func delay(cycles: int) -> void {
    let i: int = 0;
    for (i = 0; i < cycles; i++) {
        // Empty loop for delay
    }
}

// Benchmark a simple operation
func benchmark_addition() -> void {
    let iterations: int = 100000;
    let sum: int = 0;
    let i: int = 0;
    
    for (i = 0; i < iterations; i++) {
        sum += i;
    }
    
    print(sum);
    print(10);
}
\end{lstlisting}

\section{Validation and Error Checking}

\begin{lstlisting}[style=algo, caption={Input Validation}]
// Check if value is in valid range
func validate_range(value: int, min_val: int, max_val: int) -> int {
    return (value >= min_val && value <= max_val) ? 1 : 0;
}

// Check if value is even
func is_even(x: int) -> int {
    return (x % 2 == 0) ? 1 : 0;
}

// Check if value is odd
func is_odd(x: int) -> int {
    return (x % 2 != 0) ? 1 : 0;
}

// Check if value is power of 2
func is_power_of_2(x: int) -> int {
    return (x > 0 && (x & (x - 1)) == 0) ? 1 : 0;
}

// Check if number is prime
func is_prime(n: int) -> int {
    if (n <= 1) {
        return 0;
    }
    if (n <= 3) {
        return 1;
    }
    if (n % 2 == 0 || n % 3 == 0) {
        return 0;
    }
    
    let i: int = 5;
    while (i * i <= n) {
        if (n % i == 0 || n % (i + 2) == 0) {
            return 0;
        }
        i += 6;
    }
    
    return 1;
}
\end{lstlisting}

\section{Random Number Generation}

\begin{lstlisting}[style=algo, caption={Simple PRNG (Linear Congruential)}]
// Global seed (would be a global variable if supported)
func lcg_rand(seed: int) -> int {
    // LCG parameters: a=1103515245, c=12345, m=2^31
    let a: int = 1103515245;
    let c: int = 12345;
    let m: int = 2147483648;  // 2^31
    
    seed = (a * seed + c) % m;
    return seed;
}

// Generate random number in range [min, max]
func rand_range(seed: int, min_val: int, max_val: int) -> int {
    let r: int = lcg_rand(seed);
    let range: int = max_val - min_val + 1;
    return min_val + (r % range);
}

// Example usage
func random_demo() -> void {
    let seed: int = 42;
    let i: int = 0;
    
    print(82);   // 'R'
    print(97);   // 'a'
    print(110);  // 'n'
    print(100);  // 'd'
    print(111);  // 'o'
    print(109);  // 'm'
    print(58);   // ':'
    print(10);   // newline
    
    for (i = 0; i < 5; i++) {
        seed = lcg_rand(seed);
        print(seed % 100);
        print(10);
    }
}
\end{lstlisting}

% ============================================================================
% CHAPTER 27: COMPLETE ALGO PROGRAMS
% ============================================================================
\chapter{Complete Algo Programs}

\section{Classic Algorithms}

\subsection{Fibonacci Sequence}

\begin{lstlisting}[style=algo, caption={Fibonacci - Iterative}]
func fibonacci(n: int) -> int {
    if (n <= 1) {
        return n;
    }
    
    let a: int = 0;
    let b: int = 1;
    let i: int = 2;
    
    while (i <= n) {
        let temp: int = a + b;
        a = b;
        b = temp;
        i++;
    }
    
    return b;
}

func print_fibonacci_series(count: int) -> void {
    let i: int = 0;
    
    for (i = 0; i < count; i++) {
        print(fibonacci(i));
        print(32);  // space
    }
    print(10);  // newline
}

func main() -> void {
    print_fibonacci_series(15);
    // Output: 0 1 1 2 3 5 8 13 21 34 55 89 144 233 377
}
\end{lstlisting}

\subsection{Fibonacci - Recursive}

\begin{lstlisting}[style=algo, caption={Fibonacci Recursive (Inefficient)}]
func fib_recursive(n: int) -> int {
    if (n <= 1) {
        return n;
    }
    return fib_recursive(n - 1) + fib_recursive(n - 2);
}

func main() -> void {
    let i: int = 0;
    for (i = 0; i < 12; i++) {
        print(fib_recursive(i));
        print(32);
    }
    print(10);
}
\end{lstlisting}

\subsection{Factorial}

\begin{lstlisting}[style=algo, caption={Factorial - Iterative and Recursive}]
// Iterative version
func factorial_iter(n: int) -> int {
    let result: int = 1;
    let i: int = 2;
    
    for (i = 2; i <= n; i++) {
        result *= i;
    }
    
    return result;
}

// Recursive version
func factorial_rec(n: int) -> int {
    if (n <= 1) {
        return 1;
    }
    return n * factorial_rec(n - 1);
}

func main() -> void {
    let n: int = 10;
    
    print(factorial_iter(n));
    print(10);
    
    print(factorial_rec(n));
    print(10);
    // Both output: 3628800
}
\end{lstlisting}

\section{Sorting Algorithms}

\subsection{Bubble Sort}

\begin{lstlisting}[style=algo, caption={Bubble Sort Implementation}]
func bubble_sort_5(a: int, b: int, c: int, d: int, e: int) -> void {
    // Register-based sorting of 5 values
    // Demonstrated with switch-case and comparisons
    
    let temp: int = 0;
    let swapped: int = 1;
    
    while (swapped) {
        swapped = 0;
        
        // Compare a and b
        if (a > b) {
            temp = a;
            a = b;
            b = temp;
            swapped = 1;
        }
        
        // Compare b and c
        if (b > c) {
            temp = b;
            b = c;
            c = temp;
            swapped = 1;
        }
        
        // Compare c and d
        if (c > d) {
            temp = c;
            c = d;
            d = temp;
            swapped = 1;
        }
        
        // Compare d and e
        if (d > e) {
            temp = d;
            d = e;
            e = temp;
            swapped = 1;
        }
    }
    
    // Print sorted values
    print(a); print(32);
    print(b); print(32);
    print(c); print(32);
    print(d); print(32);
    print(e); print(10);
}

func main() -> void {
    bubble_sort_5(64, 34, 25, 12, 22);
    // Output: 12 22 25 34 64
}
\end{lstlisting}

\subsection{Selection Sort}

\begin{lstlisting}[style=algo, caption={Selection Sort for 5 Values}]
func selection_sort_demo() -> void {
    let v1: int = 64;
    let v2: int = 25;
    let v3: int = 12;
    let v4: int = 22;
    let v5: int = 11;
    
    let min_val: int = 0;
    let temp: int = 0;
    
    // Find minimum of all 5
    min_val = v1;
    let min_pos: int = 1;
    
    if (v2 < min_val) { min_val = v2; min_pos = 2; }
    if (v3 < min_val) { min_val = v3; min_pos = 3; }
    if (v4 < min_val) { min_val = v4; min_pos = 4; }
    if (v5 < min_val) { min_val = v5; min_pos = 5; }
    
    // Swap minimum to first position
    switch (min_pos) {
        case 2:
            temp = v1; v1 = v2; v2 = temp;
            break;
        case 3:
            temp = v1; v1 = v3; v3 = temp;
            break;
        case 4:
            temp = v1; v1 = v4; v4 = temp;
            break;
        case 5:
            temp = v1; v1 = v5; v5 = temp;
            break;
        default:
            break;
    }
    
    // Continue for remaining values...
    print(v1); print(32);
    print(v2); print(32);
    print(v3); print(32);
    print(v4); print(32);
    print(v5); print(10);
}
\end{lstlisting}

\section{Search Algorithms}

\subsection{Linear Search}

\begin{lstlisting}[style=algo, caption={Linear Search}]
func linear_search(v1: int, v2: int, v3: int, 
                   v4: int, v5: int, target: int) -> int {
    if (v1 == target) { return 0; }
    if (v2 == target) { return 1; }
    if (v3 == target) { return 2; }
    if (v4 == target) { return 3; }
    if (v5 == target) { return 4; }
    return -1;  // Not found
}

func main() -> void {
    let pos: int = linear_search(10, 20, 30, 40, 50, 30);
    
    if (pos >= 0) {
        print(70);   // 'F'
        print(111);  // 'o'
        print(117);  // 'u'
        print(110);  // 'n'
        print(100);  // 'd'
        print(32);   // space
        print(97);   // 'a'
        print(116);  // 't'
        print(32);   // space
        print(pos);
        print(10);
    }
}
\end{lstlisting}

\subsection{Binary Search}

\begin{lstlisting}[style=algo, caption={Binary Search (Conceptual)}]
func binary_search_8(v0: int, v1: int, v2: int, v3: int,
                     v4: int, v5: int, v6: int, v7: int,
                     target: int) -> int {
    // Assumes sorted array
    let low: int = 0;
    let high: int = 7;
    let mid: int = 0;
    let mid_val: int = 0;
    
    while (low <= high) {
        mid = (low + high) / 2;
        
        // Get value at mid position
        switch (mid) {
            case 0: mid_val = v0; break;
            case 1: mid_val = v1; break;
            case 2: mid_val = v2; break;
            case 3: mid_val = v3; break;
            case 4: mid_val = v4; break;
            case 5: mid_val = v5; break;
            case 6: mid_val = v6; break;
            case 7: mid_val = v7; break;
            default: mid_val = 0; break;
        }
        
        if (mid_val == target) {
            return mid;
        } else if (mid_val < target) {
            low = mid + 1;
        } else {
            high = mid - 1;
        }
    }
    
    return -1;  // Not found
}

func main() -> void {
    let result: int = binary_search_8(2, 5, 8, 12, 16, 23, 38, 56, 23);
    print(result);  // Output: 5
    print(10);
}
\end{lstlisting}

\section{Number Theory}

\subsection{Prime Number Checker}

\begin{lstlisting}[style=algo, caption={Prime Number Detection}]
func is_prime(n: int) -> int {
    if (n <= 1) {
        return 0;
    }
    if (n <= 3) {
        return 1;
    }
    if (n % 2 == 0 || n % 3 == 0) {
        return 0;
    }
    
    let i: int = 5;
    while (i * i <= n) {
        if (n % i == 0 || n % (i + 2) == 0) {
            return 0;
        }
        i += 6;
    }
    
    return 1;
}

func print_primes_up_to(limit: int) -> void {
    let n: int = 2;
    
    for (n = 2; n <= limit; n++) {
        if (is_prime(n)) {
            print(n);
            print(32);  // space
        }
    }
    print(10);  // newline
}

func main() -> void {
    print_primes_up_to(50);
    // Output: 2 3 5 7 11 13 17 19 23 29 31 37 41 43 47
}
\end{lstlisting}

\subsection{GCD and LCM}

\begin{lstlisting}[style=algo, caption={Greatest Common Divisor}]
// Euclidean algorithm
func gcd(a: int, b: int) -> int {
    while (b != 0) {
        let temp: int = b;
        b = a % b;
        a = temp;
    }
    return a;
}

func lcm(a: int, b: int) -> int {
    return (a * b) / gcd(a, b);
}

func main() -> void {
    print(gcd(48, 18));    // Output: 6
    print(10);
    
    print(lcm(12, 18));    // Output: 36
    print(10);
}
\end{lstlisting}

\section{Recursion Examples}

\subsection{Power Function}

\begin{lstlisting}[style=algo, caption={Exponentiation by Recursion}]
func power(base: int, exp: int) -> int {
    if (exp == 0) {
        return 1;
    }
    if (exp == 1) {
        return base;
    }
    
    let half: int = power(base, exp / 2);
    
    if (exp % 2 == 0) {
        return half * half;
    } else {
        return base * half * half;
    }
}

func main() -> void {
    print(power(2, 10));   // Output: 1024
    print(10);
    
    print(power(3, 5));    // Output: 243
    print(10);
}
\end{lstlisting}

\subsection{Tower of Hanoi}

\begin{lstlisting}[style=algo, caption={Tower of Hanoi}]
func hanoi(n: int, from: int, to: int, aux: int) -> void {
    if (n == 1) {
        print(77);   // 'M'
        print(111);  // 'o'
        print(118);  // 'v'
        print(101);  // 'e'
        print(32);   // space
        print(from);
        print(32);
        print(116);  // 't'
        print(111);  // 'o'
        print(32);
        print(to);
        print(10);
        return;
    }
    
    hanoi(n - 1, from, aux, to);
    
    print(77); print(111); print(118); print(101); print(32);
    print(from); print(32);
    print(116); print(111); print(32);
    print(to); print(10);
    
    hanoi(n - 1, aux, to, from);
}

func main() -> void {
    hanoi(3, 1, 3, 2);
}
\end{lstlisting}

\section{Dynamic Programming}

\subsection{Fibonacci with Memoization}

\begin{lstlisting}[style=algo, caption={Fibonacci with Register-Based Cache}]
func fib_memo(n: int) -> int {
    // Simple cache for first 10 Fibonacci numbers
    let f0: int = 0;
    let f1: int = 1;
    let f2: int = 1;
    let f3: int = 2;
    let f4: int = 3;
    let f5: int = 5;
    let f6: int = 8;
    let f7: int = 13;
    let f8: int = 21;
    let f9: int = 34;
    
    switch (n) {
        case 0: return f0;
        case 1: return f1;
        case 2: return f2;
        case 3: return f3;
        case 4: return f4;
        case 5: return f5;
        case 6: return f6;
        case 7: return f7;
        case 8: return f8;
        case 9: return f9;
        default:
            // Calculate for larger n
            let a: int = f9;
            let b: int = f8;
            let i: int = 10;
            
            while (i <= n) {
                let temp: int = a + b;
                b = a;
                a = temp;
                i++;
            }
            return a;
    }
}
\end{lstlisting}

\section{Bit Manipulation Showcase}

\subsection{Complete Bit Operations Demo}

\begin{lstlisting}[style=algo, caption={Comprehensive Bit Manipulation}]
func bit_operations_demo() -> void {
    let x: int = 42;  // Binary: 00101010
    
    print(120);  // 'x'
    print(61);   // '='
    print(x);
    print(10);
    
    // Set bit 0
    x |= 1;
    print(x);    // 43 (00101011)
    print(10);
    
    // Clear bit 1
    x &= ~2;
    print(x);    // 41 (00101001)
    print(10);
    
    // Toggle bit 3
    x ^= 8;
    print(x);    // 33 (00100001)
    print(10);
    
    // Left shift
    x <<= 1;
    print(x);    // 66 (01000010)
    print(10);
    
    // Right shift
    x >>= 2;
    print(x);    // 16 (00010000)
    print(10);
}

func main() -> void {
    bit_operations_demo();
}
\end{lstlisting}

\section{Complete Program: Calculator}

\begin{lstlisting}[style=algo, caption={Simple Integer Calculator}]
func calculator(op: int, a: int, b: int) -> int {
    let result: int = 0;
    
    switch (op) {
        case 1:  // Addition
            result = a + b;
            break;
        case 2:  // Subtraction
            result = a - b;
            break;
        case 3:  // Multiplication
            result = a * b;
            break;
        case 4:  // Division
            if (b != 0) {
                result = a / b;
            } else {
                print(69); print(114); print(114); // "Err"
                print(10);
                return 0;
            }
            break;
        case 5:  // Modulo
            if (b != 0) {
                result = a % b;
            } else {
                print(69); print(114); print(114);
                print(10);
                return 0;
            }
            break;
        case 6:  // Power
            result = 1;
            let i: int = 0;
            for (i = 0; i < b; i++) {
                result *= a;
            }
            break;
        default:
            print(73); print(110); print(118); // "Inv"
            print(10);
            return 0;
    }
    
    return result;
}

func main() -> void {
    print(calculator(1, 10, 5));   // 15
    print(10);
    
    print(calculator(3, 7, 8));    // 56
    print(10);
    
    print(calculator(6, 2, 10));   // 1024
    print(10);
}
\end{lstlisting}

\section{Complete Program: Number Guessing Game}

\begin{lstlisting}[style=algo, caption={Interactive Guessing Game Simulation}]
func guessing_game() -> void {
    let secret: int = 42;
    let guess: int = 0;
    let attempts: int = 0;
    let max_attempts: int = 7;
    
    // Simulate guesses
    let guesses: int = 0;
    
    // Pre-programmed guesses for demonstration
    guess = 50; attempts++; 
    if (guess > secret) { print(72); } // 'H'igh
    else if (guess < secret) { print(76); } // 'L'ow
    else { goto won; }
    
    guess = 25; attempts++;
    if (guess > secret) { print(72); }
    else if (guess < secret) { print(76); }
    else { goto won; }
    
    guess = 40; attempts++;
    if (guess > secret) { print(72); }
    else if (guess < secret) { print(76); }
    else { goto won; }
    
    guess = 42; attempts++;
    if (guess == secret) { goto won; }
    
    goto lost;
    
    won:
    print(87); print(105); print(110); // "Win"
    print(32);
    print(105); print(110); print(32);
    print(attempts);
    print(10);
    return;
    
    lost:
    print(76); print(111); print(115); print(116); // "Lost"
    print(10);
}

func main() -> void {
    guessing_game();
    // Output: L L L Win in 4
}
\end{lstlisting}

\section{Summary}

This chapter demonstrated complete, working Algo programs that showcase:

\begin{itemize}
    \item \textbf{All implemented features}: \code{func}, \code{let}, \code{->}, type annotations
    \item \textbf{All operators}: \code{++/--}, compound assignments (\code{+=}, \code{-=}, etc.), ternary (\code{?:}), \code{sizeof}
    \item \textbf{All control structures}: \code{if/else}, \code{while}, \code{do-while}, \code{for}, \code{switch/case}, \code{break}, \code{continue}, \code{goto}
    \item \textbf{Enums}: Compile-time constants
    \item \textbf{Register-only architecture}: No heap allocation, function parameters only
    \item \textbf{Recursion}: Full support with proper stack management
\end{itemize}

All examples in this chapter are based on the actual Algo compiler implementation and will compile and run successfully on the A1M platform.

% APPENDIX A: INSTRUCTION REFERENCE
% ============================================================================
\chapter{Complete Instruction Reference}

\section{Instruction Summary}

\begin{table}[H]
\centering
\caption{Complete A1M Instruction Set}
\small
\begin{tabularx}{\textwidth}{|l|l|l|X|}
\hline
\textbf{Opcode} & \textbf{Mnemonic} & \textbf{Format} & \textbf{Description} \\
\hline
0x00 & NOP & - & No operation \\
0x01 & HALT & - & Halt execution \\
0x02 & ADD & R & Add registers \\
0x03 & ADDI & I & Add immediate \\
0x04 & SUB & R & Subtract registers \\
0x05 & SUBI & I & Subtract immediate \\
0x06 & MUL & R & Multiply registers \\
0x07 & MULI & I & Multiply immediate \\
0x08 & DIV & R & Divide registers \\
0x09 & DIVI & I & Divide immediate \\
0x0A & MOD & R & Modulo registers \\
0x0B & MODI & I & Modulo immediate \\
0x0C & AND & R & Bitwise AND \\
0x0D & ANDI & I & AND immediate \\
0x0E & OR & R & Bitwise OR \\
0x0F & ORI & I & OR immediate \\
0x10 & XOR & R & Bitwise XOR \\
0x11 & XORI & I & XOR immediate \\
0x12 & SHL & R & Shift left \\
0x13 & SHLI & I & Shift left immediate \\
0x14 & SHR & R & Shift right logical \\
0x15 & SHRI & I & Shift right immediate \\
0x16 & SAR & R & Shift right arithmetic \\
0x17 & SARI & I & SAR immediate \\
0x20 & LOAD & M & Load 64-bit word \\
0x21 & STORE & M & Store 64-bit word \\
0x22 & LOADB & M & Load byte \\
0x23 & STOREB & M & Store byte \\
0x30 & JMP & J & Unconditional jump \\
0x31 & JZ & J & Jump if zero \\
0x32 & JNZ & J & Jump if not zero \\
0x33 & JGT & J & Jump if greater \\
0x34 & JLT & J & Jump if less \\
0x35 & JGE & J & Jump if greater/equal \\
0x36 & JLE & J & Jump if less/equal \\
0x39 & CALL & J & Function call \\
0x3A & RET & - & Return from function \\
0x40 & CMP & R & Compare registers \\
0x41 & CMPI & I & Compare immediate \\
0x42 & PUSH & R & Push to stack \\
0x43 & POP & R & Pop from stack \\
0x50 & SYSCALL & I & System call \\
\hline
\end{tabularx}
\end{table}

% ============================================================================
% APPENDIX B: REGISTER CONVENTIONS
% ============================================================================
\chapter{Register Usage Conventions}

\begin{table}[H]
\centering
\caption{A1M Calling Convention}
\begin{tabularx}{\textwidth}{|l|l|l|X|}
\hline
\textbf{Register} & \textbf{Name} & \textbf{Preserved} & \textbf{Purpose} \\
\hline
R0 & zero & N/A & Always zero \\
R1 & ret/arg1 & No & Return value / Syscall arg 1 \\
R2 & arg2 & No & Syscall argument 2 \\
R3 & sp & Yes & Stack pointer \\
R4-R11 & a0-a7 & No & Function arguments \\
R12-R30 & t0-t18 & No & Temporaries \\
R31 & ra & No & Return address \\
\hline
\end{tabularx}
\end{table}

% ============================================================================
% APPENDIX C: SYSTEM CALL REFERENCE
% ============================================================================
\chapter{System Call Reference}

\begin{table}[H]
\centering
\caption{System Call Numbers and Arguments}
\small
\begin{tabularx}{\textwidth}{|l|l|X|}
\hline
\textbf{Number} & \textbf{Name} & \textbf{Arguments (R2, R3, R4)} \\
\hline
0 & exit & exit\_code \\
1 & write & fd, buffer, length \\
2 & read & fd, buffer, length \\
3 & open & path, flags \\
4 & close & fd \\
5 & malloc & size → returns address in R1 \\
6 & free & address \\
10 & print\_int & value \\
11 & print\_str & string\_address \\
12 & read\_int & → returns value in R1 \\
13 & read\_str & buffer, max\_length \\
\hline
\end{tabularx}
\end{table}

% ============================================================================
% APPENDIX D: ALGO LANGUAGE GRAMMAR
% ============================================================================
\chapter{Algo Language Grammar}

\section{EBNF Grammar}

\begin{verbatim}
program     ::= function*

function    ::= 'fun' IDENTIFIER '(' parameters ')' 
                ('->' type)? block

parameters  ::= parameter (',' parameter)*
parameter   ::= type IDENTIFIER ('[' ']')?

block       ::= '{' statement* '}'

statement   ::= var_decl
              | assign
              | if_stmt
              | while_stmt
              | for_stmt
              | return_stmt
              | expr_stmt

var_decl    ::= type IDENTIFIER ('[' INTEGER ']')? 
                ('=' expression)? ';'

assign      ::= IDENTIFIER ('['expression']')? '=' expression ';'

if_stmt     ::= 'if' '(' expression ')' block ('else' block)?

while_stmt  ::= 'while' '(' expression ')' block

for_stmt    ::= 'for' '(' var_decl expression ';' 
                assign ')' block

return_stmt ::= 'return' expression? ';'

expression  ::= logic_or

logic_or    ::= logic_and ('||' logic_and)*
logic_and   ::= equality ('&&' equality)*
equality    ::= comparison (('=='|'!=') comparison)*
comparison  ::= term (('<'|'>'|'<='|'>=') term)*
term        ::= factor (('+'|'-') factor)*
factor      ::= unary (('*'|'/'|'%') unary)*
unary       ::= ('!'|'-') unary | primary

primary     ::= INTEGER | STRING | CHAR | 'true' | 'false'
              | IDENTIFIER
              | IDENTIFIER '(' arguments ')'
              | IDENTIFIER '[' expression ']'
              | '(' expression ')'

type        ::= 'int' | 'long' | 'char' | 'bool'
\end{verbatim}

% ============================================================================
% APPENDICES
% ============================================================================
\appendix

% ============================================================================
% APPENDIX A: COMPLETE INSTRUCTION REFERENCE
% ============================================================================
\chapter{Complete Instruction Reference}

\section{Instruction Set Summary}

This appendix provides a complete reference for all A1M instructions, including opcodes, formats, operands, and detailed descriptions.

\subsection{Instruction Format Legend}

\begin{itemize}
    \item \textbf{RD} - Destination register (5 bits, R0-R31)
    \item \textbf{RS1} - Source register 1 (5 bits, R0-R31)
    \item \textbf{RS2} - Source register 2 (5 bits, R0-R31)
    \item \textbf{IMM} - Immediate value (14 bits, signed)
    \item \textbf{ADDR} - Jump/branch address (24 bits)
    \item \textbf{OFFSET} - Memory offset (14 bits, signed)
\end{itemize}

\section{Complete Instruction List}

\begin{table}[H]
\centering
\caption{A1M Complete Instruction Set}
\label{tab:complete-instructions}
\tiny
\begin{tabularx}{\textwidth}{|l|l|l|X|l|}
\hline
\textbf{Opcode} & \textbf{Mnemonic} & \textbf{Format} & \textbf{Operation} & \textbf{Flags} \\
\hline
\multicolumn{5}{|c|}{\textbf{Control Instructions}} \\
\hline
0x00 & NOP & - & No operation & - \\
0x01 & HALT & - & Stop execution & - \\
\hline
\multicolumn{5}{|c|}{\textbf{Arithmetic Instructions}} \\
\hline
0x02 & ADD & R & RD = RS1 + RS2 & ZF, NF \\
0x03 & ADDI & I & RD = RS1 + IMM & ZF, NF \\
0x04 & SUB & R & RD = RS1 - RS2 & ZF, NF \\
0x05 & SUBI & I & RD = RS1 - IMM & ZF, NF \\
0x06 & MUL & R & RD = RS1 $\times$ RS2 & ZF, NF \\
0x07 & MULI & I & RD = RS1 $\times$ IMM & ZF, NF \\
0x08 & DIV & R & RD = RS1 / RS2 & ZF, NF \\
0x09 & DIVI & I & RD = RS1 / IMM & ZF, NF \\
0x0A & MOD & R & RD = RS1 \% RS2 & ZF, NF \\
0x0B & MODI & I & RD = RS1 \% IMM & ZF, NF \\
\hline
\multicolumn{5}{|c|}{\textbf{Logical Instructions}} \\
\hline
0x0C & AND & R & RD = RS1 \& RS2 & ZF, NF \\
0x0D & ANDI & I & RD = RS1 \& IMM & ZF, NF \\
0x0E & OR & R & RD = RS1 | RS2 & ZF, NF \\
0x0F & ORI & I & RD = RS1 | IMM & ZF, NF \\
0x10 & XOR & R & RD = RS1 $\oplus$ RS2 & ZF, NF \\
0x11 & XORI & I & RD = RS1 $\oplus$ IMM & ZF, NF \\
0x12 & SHL & R & RD = RS1 $<<$ RS2 & ZF, NF \\
0x13 & SHLI & I & RD = RS1 $<<$ IMM & ZF, NF \\
0x14 & SHR & R & RD = RS1 $>>$ RS2 (logical) & ZF, NF \\
0x15 & SHRI & I & RD = RS1 $>>$ IMM (logical) & ZF, NF \\
0x16 & SAR & R & RD = RS1 $>>$ RS2 (arithmetic) & ZF, NF \\
0x17 & SARI & I & RD = RS1 $>>$ IMM (arithmetic) & ZF, NF \\
0x18 & NOT & R & RD = $\sim$RS1 & ZF, NF \\
\hline
\multicolumn{5}{|c|}{\textbf{Memory Instructions}} \\
\hline
0x20 & LOAD & M & RD = Mem[RS1 + OFFSET] & - \\
0x21 & STORE & M & Mem[RS1 + OFFSET] = RD & - \\
0x22 & LOADB & M & RD = (byte)Mem[RS1 + OFFSET] & - \\
0x23 & STOREB & M & Mem[RS1 + OFFSET] = (byte)RD & - \\
\hline
\multicolumn{5}{|c|}{\textbf{Branch and Jump Instructions}} \\
\hline
0x30 & JMP & J & PC = ADDR & - \\
0x31 & JZ & J & if (ZF) PC = ADDR & - \\
0x32 & JNZ & J & if (!ZF) PC = ADDR & - \\
0x33 & JGT & J & if (!ZF \&\& !NF) PC = ADDR & - \\
0x34 & JLT & J & if (NF) PC = ADDR & - \\
0x35 & JGE & J & if (!NF) PC = ADDR & - \\
0x36 & JLE & J & if (ZF || NF) PC = ADDR & - \\
0x37 & JEQ & J & if (ZF) PC = ADDR & - \\
0x38 & JNE & J & if (!ZF) PC = ADDR & - \\
0x39 & CALL & J & R31 = PC + 4; PC = ADDR & - \\
0x3A & RET & - & PC = R31 & - \\
\hline
\multicolumn{5}{|c|}{\textbf{Comparison Instructions}} \\
\hline
0x40 & CMP & R & Set flags: RS1 - RS2 & ZF, NF, CF, OF \\
0x41 & CMPI & I & Set flags: RS1 - IMM & ZF, NF, CF, OF \\
\hline
\multicolumn{5}{|c|}{\textbf{Stack Instructions}} \\
\hline
0x42 & PUSH & R & SP -= 8; Mem[SP] = RS1 & - \\
0x43 & POP & R & RD = Mem[SP]; SP += 8 & - \\
\hline
\multicolumn{5}{|c|}{\textbf{System Instructions}} \\
\hline
0x50 & SYSCALL & I & System call (IMM = number) & - \\
\hline
\end{tabularx}
\end{table}

\section{Instruction Encoding Details}

\subsection{R-Type Format}

\begin{verbatim}
 31    24 23  19 18  14 13   9 8       0
┌────────┬──────┬──────┬──────┬─────────┐
│ Opcode │  RD  │ RS1  │ RS2  │ Unused  │
└────────┴──────┴──────┴──────┴─────────┘
  8 bits  5 bits 5 bits 5 bits  9 bits
\end{verbatim}

\textbf{Example:} \code{ADD R5, R10, R15}

\begin{itemize}
    \item Opcode: 0x02 (ADD)
    \item RD: 5 (00101)
    \item RS1: 10 (01010)
    \item RS2: 15 (01111)
    \item Encoding: \code{0x02 2A 78 00}
\end{itemize}

\subsection{I-Type Format}

\begin{verbatim}
 31    24 23  19 18  14 13           0
┌────────┬──────┬──────┬──────────────┐
│ Opcode │  RD  │ RS1  │  Immediate   │
└────────┴──────┴──────┴──────────────┘
  8 bits  5 bits 5 bits   14 bits
\end{verbatim}

\textbf{Example:} \code{ADDI R3, R4, 100}

\begin{itemize}
    \item Opcode: 0x03 (ADDI)
    \item RD: 3 (00011)
    \item RS1: 4 (00100)
    \item Immediate: 100 (0000000001100100)
    \item Encoding: \code{0x03 1A 06 64}
\end{itemize}

% ============================================================================
% APPENDIX B: REGISTER CONVENTIONS
% ============================================================================
\chapter{Register Usage Conventions}

\section{Register Assignment}

\begin{table}[H]
\centering
\caption{A1M Register Assignment and Usage}
\label{tab:register-assignment}
\begin{tabularx}{\textwidth}{|l|l|l|l|X|}
\hline
\textbf{Reg} & \textbf{ABI} & \textbf{Saved} & \textbf{Usage} & \textbf{Notes} \\
\hline
R0 & zero & N/A & Constant zero & Hardwired to 0, writes ignored \\
R1 & ret & No & Return value & Also syscall number \\
R2 & arg2 & No & Syscall arg 2 & Second syscall argument \\
R3 & sp & Yes & Stack pointer & Must be preserved by callee \\
R4 & a0 & No & Argument 0 & First function argument \\
R5 & a1 & No & Argument 1 & Second function argument \\
R6 & a2 & No & Argument 2 & Third function argument \\
R7 & a3 & No & Argument 3 & Fourth function argument \\
R8 & a4 & No & Argument 4 & Fifth function argument \\
R9 & a5 & No & Argument 5 & Sixth function argument \\
R10 & a6 & No & Argument 6 & Seventh function argument \\
R11 & a7 & No & Argument 7 & Eighth function argument \\
R12 & t0 & No & Temporary 0 & Caller-saved temporary \\
R13 & t1 & No & Temporary 1 & Caller-saved temporary \\
R14 & t2 & No & Temporary 2 & Caller-saved temporary \\
R15 & t3 & No & Temporary 3 & Caller-saved temporary \\
R16 & t4 & No & Temporary 4 & Caller-saved temporary \\
R17 & t5 & No & Temporary 5 & Caller-saved temporary \\
R18 & t6 & No & Temporary 6 & Caller-saved temporary \\
R19 & t7 & No & Temporary 7 & Caller-saved temporary \\
R20 & t8 & No & Temporary 8 & Caller-saved temporary \\
R21 & t9 & No & Temporary 9 & Caller-saved temporary \\
R22 & t10 & No & Temporary 10 & Caller-saved temporary \\
R23 & t11 & No & Temporary 11 & Caller-saved temporary \\
R24 & t12 & No & Temporary 12 & Caller-saved temporary \\
R25 & t13 & No & Temporary 13 & Caller-saved temporary \\
R26 & t14 & No & Temporary 14 & Caller-saved temporary \\
R27 & t15 & No & Temporary 15 & Caller-saved temporary \\
R28 & t16 & No & Temporary 16 & Caller-saved temporary \\
R29 & t17 & No & Temporary 17 & Caller-saved temporary \\
R30 & t18 & No & Temporary 18 & Caller-saved temporary \\
R31 & ra & No & Return address & Set by CALL, used by RET \\
\hline
\end{tabularx}
\end{table}

\section{Calling Convention}

\subsection{Function Prologue}

\begin{lstlisting}[style=assembly, caption={Standard Function Prologue}]
function_name:
    PUSH R31            ; Save return address
    PUSH R3             ; Save old stack pointer (optional)
    SUBI R3, R3, N      ; Allocate N bytes for locals
    ; Save callee-saved registers if used
\end{lstlisting}

\subsection{Function Epilogue}

\begin{lstlisting}[style=assembly, caption={Standard Function Epilogue}]
    ; Restore callee-saved registers
    ADDI R3, R3, N      ; Deallocate local variables
    POP R3              ; Restore stack pointer (if saved)
    POP R31             ; Restore return address
    RET                 ; Return to caller
\end{lstlisting}

\subsection{Argument Passing}

\begin{itemize}
    \item Arguments 1-8: Passed in R4-R11
    \item Additional arguments: Passed on stack
    \item Return value: R1
    \item Return address: R31 (set by CALL)
\end{itemize}

% ============================================================================
% APPENDIX C: SYSTEM CALL REFERENCE
% ============================================================================
\chapter{System Call Reference}

\section{System Call Interface}

System calls use the \code{SYSCALL} instruction with the following convention:

\begin{itemize}
    \item \textbf{R1}: System call number
    \item \textbf{R2}: First argument
    \item \textbf{R3}: Second argument  
    \item \textbf{R4}: Third argument
    \item \textbf{R1}: Return value (after syscall)
\end{itemize}

\section{System Call Table}

\begin{table}[H]
\centering
\caption{A1M System Call Numbers and Signatures}
\label{tab:syscall-reference}
\small
\begin{tabularx}{\textwidth}{|l|l|X|X|}
\hline
\textbf{Num} & \textbf{Name} & \textbf{Arguments} & \textbf{Returns} \\
\hline
0 & exit & R2: exit\_code & Does not return \\
\hline
1 & write & R2: fd, R3: buf, R4: len & R1: bytes written \\
\hline
2 & read & R2: fd, R3: buf, R4: len & R1: bytes read \\
\hline
3 & open & R2: path, R3: flags & R1: file descriptor \\
\hline
4 & close & R2: fd & R1: 0 on success \\
\hline
5 & malloc & R2: size & R1: address \\
\hline
6 & free & R2: address & R1: 0 on success \\
\hline
10 & print\_int & R2: value & - \\
\hline
11 & print\_str & R2: string\_addr & - \\
\hline
12 & read\_int & - & R1: integer value \\
\hline
13 & read\_str & R2: buffer, R3: max\_len & R1: bytes read \\
\hline
14 & print\_char & R2: character & - \\
\hline
15 & read\_char & - & R1: character \\
\hline
\end{tabularx}
\end{table}

\section{System Call Examples}

\subsection{Exit Program}

\begin{lstlisting}[style=assembly]
    MOVI R1, 0      ; syscall: exit
    MOVI R2, 0      ; exit code: 0 (success)
    SYSCALL 0
\end{lstlisting}

\subsection{Print Integer}

\begin{lstlisting}[style=assembly]
    MOVI R1, 10     ; syscall: print_int
    MOVI R2, 42     ; value to print
    SYSCALL 0
\end{lstlisting}

\subsection{Read and Write File}

\begin{lstlisting}[style=assembly]
.section .data
filename: .string "data.txt"
buffer:   .space 1024

.section .text
    ; Open file
    MOVI R1, 3          ; syscall: open
    MOVI R2, filename   ; file path
    MOVI R3, 0          ; flags: O_RDONLY
    SYSCALL 0
    MOV R10, R1         ; Save file descriptor
    
    ; Read from file
    MOVI R1, 2          ; syscall: read
    MOV R2, R10         ; file descriptor
    MOVI R3, buffer     ; buffer address
    MOVI R4, 1024       ; max bytes to read
    SYSCALL 0
    MOV R11, R1         ; Save bytes read
    
    ; Write to stdout
    MOVI R1, 1          ; syscall: write
    MOVI R2, 1          ; fd: stdout
    MOVI R3, buffer     ; buffer address
    MOV R4, R11         ; bytes to write
    SYSCALL 0
    
    ; Close file
    MOVI R1, 4          ; syscall: close
    MOV R2, R10         ; file descriptor
    SYSCALL 0
\end{lstlisting}

% ============================================================================
% APPENDIX D: ALGO LANGUAGE GRAMMAR
% ============================================================================
\chapter{Algo Language Complete Grammar}

\section{EBNF Notation}

\begin{itemize}
    \item \code{::=} - Production rule
    \item \code{|} - Alternative
    \item \code{*} - Zero or more
    \item \code{+} - One or more
    \item \code{?} - Optional (zero or one)
    \item \code{( )} - Grouping
    \item \code{UPPERCASE} - Terminal (keyword or token type)
    \item \code{lowercase} - Non-terminal
\end{itemize}

\section{Complete Grammar}

\begin{verbatim}
program         ::= declaration*

declaration     ::= function_decl
                  | var_decl

function_decl   ::= 'fun' IDENTIFIER '(' parameter_list? ')' 
                    ('->' type)? block

parameter_list  ::= parameter (',' parameter)*

parameter       ::= type IDENTIFIER ('[' ']')?

var_decl        ::= type IDENTIFIER ('[' INTEGER ']')? 
                    ('=' expression)? ';'

block           ::= '{' statement* '}'

statement       ::= var_decl
                  | expr_stmt
                  | if_stmt
                  | while_stmt
                  | for_stmt
                  | return_stmt
                  | block

expr_stmt       ::= expression ';'

if_stmt         ::= 'if' '(' expression ')' statement
                    ('else' statement)?

while_stmt      ::= 'while' '(' expression ')' statement

for_stmt        ::= 'for' '(' (var_decl | expr_stmt | ';')
                    expression? ';'
                    expression? ')' statement

return_stmt     ::= 'return' expression? ';'

expression      ::= assignment

assignment      ::= (IDENTIFIER ('[' expression ']')?)
                    '=' assignment
                  | logic_or

logic_or        ::= logic_and ('||' logic_and)*

logic_and       ::= equality ('&&' equality)*

equality        ::= comparison (('==' | '!=') comparison)*

comparison      ::= term (('<' | '>' | '<=' | '>=') term)*

term            ::= factor (('+' | '-') factor)*

factor          ::= unary (('*' | '/' | '%') unary)*

unary           ::= ('!' | '-') unary
                  | postfix

postfix         ::= primary ('[' expression ']'
                  | '(' argument_list? ')')*

primary         ::= INTEGER
                  | STRING
                  | CHAR
                  | 'true'
                  | 'false'
                  | IDENTIFIER
                  | '(' expression ')'

argument_list   ::= expression (',' expression)*

type            ::= 'int' | 'long' | 'char' | 'bool'
\end{verbatim}

\section{Operator Precedence}

\begin{table}[H]
\centering
\caption{Algo Operator Precedence (Highest to Lowest)}
\label{tab:operator-precedence}
\begin{tabularx}{\textwidth}{|l|X|l|}
\hline
\textbf{Precedence} & \textbf{Operators} & \textbf{Associativity} \\
\hline
1 (Highest) & Function call \code{()}, Array access \code{[]} & Left \\
\hline
2 & Unary \code{!}, \code{-} & Right \\
\hline
3 & Multiplicative \code{*}, \code{/}, \code{\%} & Left \\
\hline
4 & Additive \code{+}, \code{-} & Left \\
\hline
5 & Relational \code{<}, \code{>}, \code{<=}, \code{>=} & Left \\
\hline
6 & Equality \code{==}, \code{!=} & Left \\
\hline
7 & Logical AND \code{\&\&} & Left \\
\hline
8 & Logical OR \code{||} & Left \\
\hline
9 (Lowest) & Assignment \code{=} & Right \\
\hline
\end{tabularx}
\end{table}

\section{Type System}

\subsection{Primitive Types}

\begin{table}[H]
\centering
\caption{Algo Primitive Types}
\begin{tabularx}{\textwidth}{|l|l|l|X|}
\hline
\textbf{Type} & \textbf{Size} & \textbf{Range} & \textbf{Description} \\
\hline
int & 64 bits & $-2^{63}$ to $2^{63}-1$ & Signed integer \\
\hline
long & 64 bits & $-2^{63}$ to $2^{63}-1$ & Alias for int \\
\hline
char & 8 bits & 0 to 255 & Single character (ASCII) \\
\hline
bool & 8 bits & 0 or 1 & Boolean (true/false) \\
\hline
\end{tabularx}
\end{table}

\subsection{Array Types}

Arrays are declared with fixed sizes:

\begin{verbatim}
int arr[10];        // Array of 10 integers
char str[256];      // Character array (string)
bool flags[8];      // Array of booleans
\end{verbatim}

% ============================================================================
% APPENDIX E: ERROR MESSAGES
% ============================================================================
\chapter{Compiler Error Messages}

\section{Lexical Errors}

\begin{table}[H]
\centering
\caption{Lexer Error Messages}
\begin{tabularx}{\textwidth}{|l|X|}
\hline
\textbf{Error Code} & \textbf{Message} \\
\hline
LEX001 & Unterminated string literal \\
LEX002 & Invalid character in input \\
LEX003 & Integer literal too large \\
LEX004 & Invalid escape sequence in string \\
LEX005 & Unterminated comment \\
\hline
\end{tabularx}
\end{table}

\section{Syntax Errors}

\begin{table}[H]
\centering
\caption{Parser Error Messages}
\begin{tabularx}{\textwidth}{|l|X|}
\hline
\textbf{Error Code} & \textbf{Message} \\
\hline
SYN001 & Expected token X, found Y \\
SYN002 & Unexpected end of file \\
SYN003 & Missing semicolon \\
SYN004 & Unmatched parenthesis/bracket/brace \\
SYN005 & Invalid function declaration \\
SYN006 & Invalid expression \\
\hline
\end{tabularx}
\end{table}

\section{Semantic Errors}

\begin{table}[H]
\centering
\caption{Semantic Error Messages}
\begin{tabularx}{\textwidth}{|l|X|}
\hline
\textbf{Error Code} & \textbf{Message} \\
\hline
SEM001 & Undefined variable: X \\
SEM002 & Duplicate variable declaration: X \\
SEM003 & Type mismatch: expected X, got Y \\
SEM004 & Function X not defined \\
SEM005 & Argument count mismatch for function X \\
SEM006 & Cannot assign to non-lvalue \\
SEM007 & Array index must be integer \\
SEM008 & Return type mismatch \\
\hline
\end{tabularx}
\end{table}

% ============================================================================
% BIBLIOGRAPHY
% ============================================================================
\begin{thebibliography}{99}

\bibitem{hennessy2017}
J. L. Hennessy and D. A. Patterson,
\textit{Computer Architecture: A Quantitative Approach},
6th ed. Morgan Kaufmann, 2017.

\bibitem{patterson2013}
D. A. Patterson and J. L. Hennessy,
\textit{Computer Organization and Design: The Hardware/Software Interface},
5th ed. Morgan Kaufmann, 2013.

\bibitem{aho2006}
A. V. Aho, M. S. Lam, R. Sethi, and J. D. Ullman,
\textit{Compilers: Principles, Techniques, and Tools},
2nd ed. Addison-Wesley, 2006.

\bibitem{muchnick1997}
S. S. Muchnick,
\textit{Advanced Compiler Design and Implementation},
Morgan Kaufmann, 1997.

\bibitem{waterman2016}
A. Waterman and K. Asanović, Eds.,
\textit{The RISC-V Instruction Set Manual, Volume I: User-Level ISA},
RISC-V Foundation, 2016.

\bibitem{sweetman2006}
D. Sweetman,
\textit{See MIPS Run},
2nd ed. Morgan Kaufmann, 2006.

\bibitem{thomas2009}
D. E. Thomas and P. R. Moorby,
\textit{The Verilog Hardware Description Language},
5th ed. Springer, 2009.

\bibitem{harris2012}
S. L. Harris and D. M. Harris,
\textit{Digital Design and Computer Architecture},
2nd ed. Morgan Kaufmann, 2012.

\end{thebibliography}

\end{document}
